% 04 — FRAGMENTO SUCESOR HODGE Y BIRCH--SWINNERTON-DYER
% =======================================================
% Inserción autónoma prevista tras el marcador de integración del Libro IV.
% No importa fuentes probatorios ni modifica las tres realizaciones
% precedentes.
%
% Trazabilidad técnica no narrativa:
%   Hodge: CHI-II-HDG-01/02/03, CHI-IV-HDG-01/C01.
%   BSD:   CHI-II-BSD-01/02/03, CHI-IV-BSD-01/02/C01.
%   Fuentes: celda APP--Mayer--Vietoris, totalizador coend y controles
%   elípticos; producto fibrado genealógico de Manin, publicación acoplada
%   Kato--PT, dualidad ortogonal reducida, relleno finita--singular y regulador
%   Bockstein derivado. Los hashes y estados de ejecución permanecen en el
%   expediente técnico de la edición sucesora.

\ifdefined\HMTSucesoraStructureTrace
  \typeout{HMT-SUCCESSOR 04: Hodge and Birch--Swinnerton-Dyer realizations}
\fi

\chapter{Completeness of enriched cohomological history and algebraicity of rational Hodge classes}
\chaptermark{Cohomological completeness and Hodge algebraicity}

\subsection*{Object, domain and starting data}

Let \(X\) be a smooth projective variety over \(\mathbb C\) and let
\(p\geq0\). We write
\begin{equation}\label{eq:hodge-suc-dominio}
 \operatorname{Hdg}^{p}(X)_{\mathbb Q}
 :=H^{2p}(X,\mathbb Q)
   \cap H^{p,p}(X,\mathbb C).
\end{equation}
Let \(\mathbf{Perf}^{\geq p}(X)\) be the idempotent-complete category of
perfect complexes whose support has codimension at least \(p\). In this
section the abbreviation is used
\[
 \operatorname{Perf}^{\geq p}(X)_{\mathbb Q}
 :=K_0\!\left(\mathbf{Perf}^{\geq p}(X)\right)
   \otimes_{\mathbb Z}\mathbb Q;
\]
thus, \(R_{\mathrm{Hdg}}\) publishes the class of the constructed perfect
complex. The starting point is not a bare class of
\eqref{eq:hodge-suc-dominio}, but the limit of surviving histories
\begin{equation}\label{eq:hodge-suc-xchi}
 X_{\chi}(X,p)
 =\varprojlim_{K}\operatorname{Surv}_{K}(X,p).
\end{equation}
Each element of \(X_{\chi}(X,p)\) preserves sheet, orientation, quotient,
carry, boundary, route and memory. In particular, the return of the joint
nonadic connection satisfies
\begin{equation}\label{eq:hodge-suc-retorno}
 \Gamma _9=\operatorname{Hol}_{C_9}(\nabla_{\mathrm{TPK}}),
 \qquad g(k+9)=g(k),
 \qquad \widehat x(k+9)\neq\widehat x(k)
 \quad\text{in general}.
\end{equation}
The equality preserves the phase; the inequality preserves the depth of the
history. This distinction rules out interpreting a return as a
reinitialization.

Route composition is governed by the carry cocycle
\begin{equation}\label{eq:hodge-suc-carry}
 \operatorname{Carr}(\gamma _2\gamma _1)
 =\operatorname{Carr}(\gamma _2)H(\gamma _1)
  +Y(\gamma _2)\operatorname{Carr}(\gamma _1),
\end{equation}
and accumulated memory by
\begin{equation}\label{eq:hodge-suc-memoria}
 S_0=\sum_{r<N}(M_9^{*})^{r}D_9^{*}D_9M_9^{r}
     +(M_9^{*})^{N}S_0M_9^{N}.
\end{equation}
Equations \eqref{eq:hodge-suc-xchi}--\eqref{eq:hodge-suc-memoria}
are data constructed in Part II. They are specialized here; they are not
reconstructed from the cohomological codomain.

\subsection*{Local APP--Mayer--Vietoris cell}

Let \(Z,W\subset X\) be closed incidences with a
Tor-independent intersection and oriented supports of the declared codimension. For
integers \(a,b\geq1\), define
\begin{equation}\label{eq:hodge-suc-kappa}
 \begin{aligned}
 \kappa_{a,b}:\mathcal O_Z^{a}\oplus\mathcal O_W^{b}
 &\longrightarrow \mathcal O_{Z\cap W}^{ab},\\
 ((u_i)_{i=1}^{a},(v_j)_{j=1}^{b})
 &\longmapsto
 \bigl(u_i|_{Z\cap W}-v_j|_{Z\cap W}\bigr)_{i,j}.
 \end{aligned}
\end{equation}
On a fiber of \(Z\cap W\), the diagonal kernel of
\(\kappa_{a,b}\) has rank one and its cokernel has rank
\begin{equation}\label{eq:hodge-suc-defecto-rango}
 ab-(a+b-1)=(a-1)(b-1).
\end{equation}
If \((\rho_{\Sigma},c_{\Sigma})\) and
\((\rho_{\Pi},c_{\Pi})\) are, respectively, the residues and carries of
the additive and multiplicative sheets of APP, the same quantity is expressed by
\begin{equation}\label{eq:hodge-suc-celda-carry}
 (a-1)(b-1)
 =\rho_{\Pi}-\rho_{\Sigma}+1
  +9\bigl(c_{\Pi}-c_{\Sigma}\bigr).
\end{equation}
Thus, the local defect is not an external correction: it is the publication in
\(\operatorname{Perf}\) of the difference between the two sheets with its carry.

Let \(\tau\) be the exchange of the two summands of the source and let \(P\)
be the transposition of the target indices. The TRIT orientation fixes
\begin{equation}\label{eq:hodge-suc-orientacion-kappa}
 \kappa_{b,a}\tau=-P\kappa_{a,b}.
\end{equation}
The sign of \eqref{eq:hodge-suc-orientacion-kappa} is part of the oriented
object. Omitting it would identify opposite routes and destroy the
naturality of the localized character.

\begin{proposition}[Perfect realization of the local cell]
\label{prop:hodge-suc-celda-perf}
The cone
\begin{equation}\label{eq:hodge-suc-cone-kappa}
 \mathcal C_{a,b}(Z,W):=\operatorname{Cone}(\kappa_{a,b})
\end{equation}
is a perfect complex with the prescribed support. Its localized class
preserves the rank defect \eqref{eq:hodge-suc-defecto-rango}, the carry
\eqref{eq:hodge-suc-celda-carry} and the orientation
\eqref{eq:hodge-suc-orientacion-kappa}.
\end{proposition}

\begin{proof}
Tor-independence identifies the derived restriction to the intersection with the
ordinary restriction used in \eqref{eq:hodge-suc-kappa}. The sources and
target are perfect on their supports; by stability of
\(\operatorname{Perf}\) under cones, so is
\(\mathcal C_{a,b}(Z,W)\). The fiber-by-fiber calculation gives
\eqref{eq:hodge-suc-defecto-rango}; APP division of the sum and
product data gives \eqref{eq:hodge-suc-celda-carry}. Finally, exchanging
\(Z,a\) with \(W,b\) changes each difference to its opposite, proving
\eqref{eq:hodge-suc-orientacion-kappa}. Since the localized Chern character
is additive in distinguished triangles, the three pieces of information pass to the
class of the cone.
\end{proof}

\subsection*{Totalization of histories and preservation of the terminal}

Let \(J_N\) be the finite category of compatible histories at depth
\(N\), let \(\Gamma\) be the diagram of incidence complexes obtained from
\eqref{eq:hodge-suc-cone-kappa} and let \(\mathsf W\) be the right module that
records route orientations and weights. The derived totalization is the
coend
\begin{equation}\label{eq:hodge-suc-coend}
 K_{X,p,N}
 :=\mathsf W\otimes^{\mathbf L}_{\mathbb Z[J_N]}\Gamma.
\end{equation}
The survival links form a compatible system in \(N\). The
holonomy \(\Gamma_9\) induces an endomorphism
\(U_{\Gamma_9}\) of the totalizer and the terminal object is preserved as
\begin{equation}\label{eq:hodge-suc-terminal}
 \mathcal T_{X,p}
 :=\operatorname{Cone}\bigl(I-U_{\Gamma_9}\bigr).
\end{equation}
The cone \eqref{eq:hodge-suc-terminal} prevents totalization from confusing
a repeated phase with the same history: the visible value may return,
but the terminal difference retains the memory displacement.

Compatibility between \eqref{eq:hodge-suc-coend} and the limit
\eqref{eq:hodge-suc-xchi} defines the realization morphism
\begin{equation}\label{eq:hodge-suc-realizacion}
 R_{\mathrm{Hdg}}:
 X_{\chi}(X,p)\longrightarrow
 \operatorname{Perf}^{\geq p}(X)_{\mathbb Q}.
\end{equation}
This morphism is fixed at the source: it does not receive a Hodge class
or a representative it must recognize as input.

\begin{theorem}[Hodge publication identity]
\label{thm:hodge-suc-identidad}
On \(X_{\chi}(X,p)\) the identity commutes
\begin{equation}\label{eq:hodge-suc-identidad}
 \operatorname{cl}_{X,p}\circ
 \operatorname{ch}^{\mathrm{loc}}_{p}\circ
 R_{\mathrm{Hdg}}
 =\operatorname{ev}^{\mathrm{Hdg}}_{\infty,X,p}.
\end{equation}
Here
\(\operatorname{ch}^{\mathrm{loc}}_{p}\) takes the codimension
\(p\) component of the localized character and
\(\operatorname{cl}_{X,p}\) is its cohomological publication.
\end{theorem}

\begin{proof}
In a cell, the equality is the compatibility of the localized character with
the restriction, the Mayer--Vietoris triangle and the cohomological class. The
orientation is fixed by \eqref{eq:hodge-suc-orientacion-kappa};
therefore, no sign remains to be chosen. Additivity of the Chern character
transports the equality to each cone \eqref{eq:hodge-suc-cone-kappa}.
Naturality with respect to the arrows of \(J_N\) makes it descend to the coend
\eqref{eq:hodge-suc-coend}. The carry law
\eqref{eq:hodge-suc-carry} guarantees its compatibility under composition and
\eqref{eq:hodge-suc-terminal} preserves the terminal term. Passing to the
projective limit yields two natural transformations with equal
components on every finite cylinder; hence they coincide and
\eqref{eq:hodge-suc-identidad} follows.
\end{proof}

The completeness theorem of the atlas of enriched histories, applied to
this specialization, establishes
\begin{equation}\label{eq:hodge-suc-completitud}
 \operatorname{ev}^{\mathrm{Hdg}}_{\infty,X,p}
 \bigl(X_{\chi}(X,p)\bigr)
 =\operatorname{Hdg}^{p}(X)_{\mathbb Q}.
\end{equation}
This is the explicit starting structure of the Hodge specialization.

% Ampliación sustantiva de la completitud coinductiva Hodge.
% Módulo destinado al capítulo Hodge de la Parte IV.

\section{Coinductive unfolding of Hodge completeness}
\label{sec:amp-hodge-completitud-desplegada}

The completeness equality
\eqref{eq:hodge-suc-completitud} can be unfolded without starting from an algebraic
class or introducing into the source the conclusion one wishes to publish. The
object being completed is the system of enriched histories; the Hodge
class intervenes only as compatible evaluation data in its
finite truncations.

\subsection{Finite observation systems and compatible fibers}
\label{subsec:amp-hodge-sistemas-finitos}

For each depth \(N\geq0\), let \(J_N(X,p)\) be the finite category of
incidence charts, oriented routes and memory terminals of depth at
most \(N\). We write
\begin{equation}\label{eq:amp-hodge-surv-n}
 \mathscr S_N(X,p):=\operatorname{Surv}_{J_N}(X,p),
 \qquad
 \rho_{N+1,N}:\mathscr S_{N+1}(X,p)\longrightarrow\mathscr S_N(X,p)
\end{equation}
for the rationalized space of surviving states and its restriction. The
finite readers take values in the module of cylinder observations
\(\mathscr H_N^p(X)\):
\begin{equation}\label{eq:amp-hodge-evaluacion-finita}
 e_N^{\mathrm{Hdg}}:\mathscr S_N(X,p)\longrightarrow
 \mathscr H_N^p(X),
 \qquad
 q_{N+1,N}:\mathscr H_{N+1}^p(X)\longrightarrow\mathscr H_N^p(X).
\end{equation}
Here \(\mathscr H_N^p(X)\) records only the evaluations on the incidences
present in \(J_N\), but preserves their rational coefficients and
orientation. If
\(q_N:\operatorname{Hdg}^p(X)_{\mathbb Q}\to\mathscr H_N^p(X)\) is the
restriction to the same finite set of observations, we have
\begin{equation}\label{eq:amp-hodge-cuadrado-restriccion}
 e_N^{\mathrm{Hdg}}\rho_{N+1,N}
 =q_{N+1,N}e_{N+1}^{\mathrm{Hdg}},
 \qquad
 q_{N+1,N}q_{N+1}=q_N.
\end{equation}

For \(\alpha\in\operatorname{Hdg}^p(X)_{\mathbb Q}\), define the compatible finite
fiber
\begin{equation}\label{eq:amp-hodge-fibra-alpha}
 \mathscr F_N(\alpha):=
 \left\{s_N\in\mathscr S_N(X,p):
 e_N^{\mathrm{Hdg}}(s_N)=q_N(\alpha)\right\}.
\end{equation}
The word \emph{finite} refers to the depth and the incidence
diagram, not to a loss of rational coefficients or to a selection
of a representative from \(X\) alone.

\begin{proposition}[Finite extension with memory]
\label{prop:amp-hodge-extension-finita}
For every \(\alpha\) as in \eqref{eq:amp-hodge-fibra-alpha}, the fibers
\(\mathscr F_N(\alpha)\) are nonempty and the restrictions induce surjective
maps
\begin{equation}\label{eq:amp-hodge-mittag-leffler}
 \rho_{N+1,N}:\mathscr F_{N+1}(\alpha)\twoheadrightarrow
 \mathscr F_N(\alpha).
\end{equation}
The extension preserves sheet, TRIT sign, quotient, carry, support, boundary and
memory terminal.
\end{proposition}

\begin{proof}
The assertion is the Hodge specialization of the extension theorem for the atlas
constructed in Part II. Its mechanism should be made visible. Let
\(s_N\in\mathscr F_N(\alpha)\) and choose a prolongation of charts
\(\widetilde s_{N+1}\) preserving the history of \(s_N\). Its discrepancy at
the new level is
\begin{equation}\label{eq:amp-hodge-discrepancia}
 \delta_{N+1}:=q_{N+1}(\alpha)
 -e_{N+1}^{\mathrm{Hdg}}(\widetilde s_{N+1}),
 \qquad q_{N+1,N}(\delta_{N+1})=0,
\end{equation}
where the last equality comes from
\eqref{eq:amp-hodge-cuadrado-restriccion}. The
APP--Mayer--Vietoris cells publish exactly that relative kernel through the
cones of the morphisms \(\kappa_{a,b}\); their defect
\((a-1)(b-1)\) is lifted with the residue and carry of
\eqref{eq:hodge-suc-celda-carry}. By the finite extension theorem there is a
relative correction \(\eta_{N+1}(\delta_{N+1})\), supported on the new
charts, such that
\begin{equation}\label{eq:amp-hodge-extension-operativa}
 \rho_{N+1,N}\eta_{N+1}(\delta_{N+1})=0,
 \qquad
 e_{N+1}^{\mathrm{Hdg}}\eta_{N+1}(\delta_{N+1})=\delta_{N+1}.
\end{equation}
Therefore
\begin{equation}\label{eq:amp-hodge-ext}
 \operatorname{Ext}_{N+1}(s_N,\alpha):=
 \widetilde s_{N+1}+\eta_{N+1}(\delta_{N+1})
 \in\mathscr F_{N+1}(\alpha)
\end{equation}
restricts to \(s_N\). The orientation
\(\kappa_{b,a}\tau=-P\kappa_{a,b}\) fixes the sign, the cocycle law
\eqref{eq:hodge-suc-carry} fixes the composition and
\(\operatorname{Cone}(I-U_{\Gamma_9})\) preserves the terminal. The initial level
is obtained from the based seed atlas. Induction proves nonemptiness and
surjectivity at all levels.
\end{proof}

\subsection{Passage to the limit and order of construction}
\label{subsec:amp-hodge-paso-limite}

\begin{theorem}[Unfolded coinductive completeness]
\label{thm:amp-hodge-completitud-coinductiva}
With the data already constructed in Part II---the systems
\eqref{eq:amp-hodge-surv-n}--\eqref{eq:amp-hodge-evaluacion-finita}, the
compatibility \eqref{eq:amp-hodge-cuadrado-restriccion}, the finite extension
of proposition \ref{prop:amp-hodge-extension-finita}, the separation of
cylinder evaluations and preservation of the terminal---for each
\(\alpha\in\operatorname{Hdg}^p(X)_{\mathbb Q}\), there is an enriched
history
\begin{equation}\label{eq:amp-hodge-xi-limite}
 \xi_\alpha=(s_0,s_1,s_2,\ldots)
 \in\varprojlim_N\mathscr F_N(\alpha)
 \subset X_\chi(X,p)
\end{equation}
which satisfies
\begin{equation}\label{eq:amp-hodge-evaluacion-alpha}
 \operatorname{ev}^{\mathrm{Hdg}}_{\infty,X,p}(\xi_\alpha)=\alpha.
\end{equation}
In particular, completeness is obtained before applying the perfect
realizer and before forming any Chow class.
\end{theorem}

\begin{proof}
Choose \(s_0\in\mathscr F_0(\alpha)\). Once \(s_N\) is constructed, the
surjectivity \eqref{eq:amp-hodge-mittag-leffler} allows us to choose
\(s_{N+1}\in\mathscr F_{N+1}(\alpha)\) with
\(\rho_{N+1,N}s_{N+1}=s_N\). Thus one obtains
\eqref{eq:amp-hodge-xi-limite}. Equivalently, the system of fibers
satisfies the Mittag--Leffler condition and produces no
\(\varprojlim^1\) obstruction term. By the definition of each fiber,
\begin{equation}\label{eq:amp-hodge-evaluaciones-xi}
 e_N^{\mathrm{Hdg}}(s_N)=q_N(\alpha)
 \qquad(N\geq0).
\end{equation}
The separation of cylinder observations identifies the unique element
of the limit of the \(q_N(\alpha)\) with \(\alpha\), which proves
\eqref{eq:amp-hodge-evaluacion-alpha}. The terminal cone prevents replacing the
sequence \((s_N)_N\) with the repetition of a visible root; therefore
\(\xi_\alpha\) preserves the rigidification used in passing to the limit.
\end{proof}

The finite publication is arranged in commutative squares
\begin{equation}\label{eq:amp-hodge-cuadrado-finito}
 \begin{array}{ccc}
 \mathscr S_N(X,p)&\xrightarrow{\ R_{\mathrm{Hdg},N}\ }&
 K_0(\mathbf{Perf}^{\geq p}(X))\otimes\mathbb Q\\[1mm]
 \big\downarrow\scriptstyle e_N^{\mathrm{Hdg}}&&
 \big\downarrow\scriptstyle
   \operatorname{cl}_{X,p}\circ\operatorname{ch}^{\mathrm{loc}}_p\\[1mm]
 \mathscr H_N^p(X)&\xrightarrow{\ \operatorname{pub}_N\ }&
 H^{2p}(X,\mathbb Q).
 \end{array}
\end{equation}
The naturality of these squares with respect to \(\rho_{N+1,N}\) gives, in the
limit, the identity \eqref{eq:hodge-suc-identidad}. The causal order is,
therefore,
\begin{equation}\label{eq:amp-hodge-orden-causal}
 \alpha\longmapsto(q_N\alpha)_N
 \longmapsto\xi_\alpha
 \longmapsto R_{\mathrm{Hdg}}(\xi_\alpha)
 \longmapsto
 \beta_\alpha:=\operatorname{ch}^{\mathrm{loc}}_p
   R_{\mathrm{Hdg}}(\xi_\alpha).
\end{equation}
Only in the last step does \(\beta_\alpha\) appear. Applying the limit square
to \eqref{eq:amp-hodge-evaluacion-alpha} yields
\begin{equation}\label{eq:amp-hodge-publicacion-final}
 \operatorname{cl}_{X,p}(\beta_\alpha)=
 \operatorname{ev}^{\mathrm{Hdg}}_{\infty,X,p}(\xi_\alpha)=\alpha.
\end{equation}

\subsection{Local check on the quadric
\texorpdfstring{\(Q^4\)}{Q4}}
\label{subsec:amp-hodge-q4}

The split quadric
\begin{equation}\label{eq:amp-hodge-q4}
 Q^4=\{x_0x_3+x_1x_4+x_2x_5=0\}\subset\mathbb P^5_{\mathbb C}
\end{equation}
has two based families of maximal planes, with representatives
\(\Pi_+\) and \(\Pi_-\). If
\(a=\operatorname{cl}[\Pi_+]\) and
\(b=\operatorname{cl}[\Pi_-]\), then
\begin{equation}\label{eq:amp-hodge-q4-bases}
 \operatorname{CH}^2(Q^4)=\mathbb Za\oplus\mathbb Zb,
 \qquad
 H^4(Q^4,\mathbb Z(2))\cap H^{2,2}(Q^4)
 =\mathbb Za\oplus\mathbb Zb.
\end{equation}
The based map
\begin{equation}\label{eq:amp-hodge-q4-beta}
 \beta_{Q^4,2}(ra+sb):=r[\Pi_+]+s[\Pi_-]
\end{equation}
satisfies exactly
\begin{equation}\label{eq:amp-hodge-q4-beta-identidades}
 \partial_{\mathrm{mot}}\beta_{Q^4,2}=0,
 \qquad
 \operatorname{cl}_{Q^4}\beta_{Q^4,2}=I.
\end{equation}

For the six ports, let
\begin{equation}\label{eq:amp-hodge-q4-matriz}
 A=\begin{pmatrix}
 2&2&2&1&2&1\\
 2&1&2&2&1&1\\
 1&0&0&1&0&2\\
 0&0&1&1&1&0\\
 2&2&2&0&2&1\\
 2&1&2&1&0&0
 \end{pmatrix},
 \qquad \det A=1,
 \qquad \mathbb A=A\otimes I.
\end{equation}
The division \(x_k\mathbb A=u_k+3x_{k+1}\) and the reduction of
\eqref{eq:amp-hodge-q4-beta} define, gate by gate,
\begin{equation}\label{eq:amp-hodge-q4-uk}
 U_{k,j}:=\overline\beta_{Q^4,2}(u_{k,j}),
 \qquad
 D_{\mathrm{mot}}U_k=0,
 \qquad
 H^0(R_{\mathrm{mot}}\otimes\mathbb F_3)(U_k)=u_k.
\end{equation}
Iterating nine times yields the telescoping identity
\begin{equation}\label{eq:amp-hodge-q4-telescopia}
 x_0\mathbb A^9
 =R_{\mathrm{mot}}\!\left(
   \sum_{k=0}^{8}3^kU_k\mathbb A^{8-k}\right)+3^9x_9.
\end{equation}
Equations \eqref{eq:amp-hodge-q4-beta-identidades}--
\eqref{eq:amp-hodge-q4-telescopia} verify local closure, port compatibility
and nonadic memory in a rank-two model. They constitute a
specialized validator of the extension mechanism; general completeness
comes from the coinductive system of the preceding theorems, not from an
extrapolation from \(Q^4\).

\subsection{Premises and refutation criteria}
\label{subsec:amp-hodge-falsadores}

The proof uses exactly five data: the finite atlas of based
charts; the relative extension of proposition
\ref{prop:amp-hodge-extension-finita}; the restriction compatibility
\eqref{eq:amp-hodge-cuadrado-restriccion}; the separation of cylinder
evaluations; and the commutativity of the publication
\eqref{eq:amp-hodge-cuadrado-finito}. Each has a locatable refutation
criterion:
\begin{enumerate}
 \item an empty fiber \(\mathscr F_N(\alpha)\) would refute finite coverage;
 \item failure of \eqref{eq:amp-hodge-mittag-leffler} would prevent constructing the
       compatible history;
 \item two distinct classes with all truncations equal would refute the
       separation of the reader;
 \item a noncommutative finite square would refute
       \eqref{eq:amp-hodge-publicacion-final};
 \item replacing the terminal with a memoryless root would invalidate the
       compatibility of the sequence, even if it preserved its visible value.
\end{enumerate}
These criteria target specific arrows of the argument. The obstruction to a
finite selector depending only on \(X\) affects none of them, because
the construction preserves the based antecedent \(\xi_\alpha\).

\begin{theorem}[Generation of the algebraic class]
\label{thm:hodge-suc-generacion}
For every \(\alpha\in\operatorname{Hdg}^{p}(X)_{\mathbb Q}\) there exists
\(\xi_{\alpha}\in X_{\chi}(X,p)\) such that, on defining
\begin{equation}\label{eq:hodge-suc-beta}
 \beta_{\alpha}:=
 \operatorname{ch}^{\mathrm{loc}}_{p}
 \bigl(R_{\mathrm{Hdg}}(\xi_{\alpha})\bigr)
 \in\operatorname{CH}^{p}(X)_{\mathbb Q},
\end{equation}
we have
\begin{equation}\label{eq:hodge-suc-cl-beta}
 \operatorname{cl}_{X,p}(\beta_{\alpha})=\alpha.
\end{equation}
\end{theorem}

\begin{proof}
By the completeness theorem
\eqref{eq:hodge-suc-completitud}, every \(\alpha\) has an antecedent
\(\xi_{\alpha}\) before applying the perfect reader. One then forms
\(\beta_{\alpha}\) through \eqref{eq:hodge-suc-beta}. The identity
\eqref{eq:hodge-suc-identidad} gives
\[
 \operatorname{cl}_{X,p}(\beta_{\alpha})
 =\operatorname{ev}^{\mathrm{Hdg}}_{\infty,X,p}(\xi_{\alpha})
 =\alpha.
\]
The rigidification lies in \(\xi_{\alpha}\), where base,
orientation, boundary and route are preserved. Therefore, choosing an antecedent does not
become a natural selector depending only on \(X\), nor has the
sought surjectivity been introduced as a premise of the map
\(R_{\mathrm{Hdg}}\).
\end{proof}

\subsection*{Hypotheses and scope}

Theorem \ref{thm:hodge-suc-generacion} uses exactly: the smooth projective
variety \(X\) and the integer \(p\); the complete atlas
\(X_{\chi}(X,p)\) and its completeness theorem
\eqref{eq:hodge-suc-completitud}; the orientations and supports of its charts; the
Tor-independence of the local intersections; the localized Chern character and
the class map defined with the same convention; and preservation of the
terminal \eqref{eq:hodge-suc-terminal}. With these explicit initial
data, \eqref{eq:hodge-suc-cl-beta} is an exact result for the entire
domain \eqref{eq:hodge-suc-dominio}.

The double publication and preservation of history are achieved through the
cell \(\kappa_{a,b}\), its cone, totalization by coend and the map
\(R_{\mathrm{Hdg}}\) with values in \(\operatorname{Perf}\). The identity
\eqref{eq:hodge-suc-identidad} and the generation
\eqref{eq:hodge-suc-cl-beta} are the results of the specialization; their
focal verifiers check support, rank, sign, carry, naturality,
terminal and commutativity of the publication.

\subsection{Final mathematical recognition of Hodge}

The final translation of \eqref{eq:hodge-suc-cl-beta} is
\begin{equation}\label{eq:hodge-suc-reconocimiento-final}
 \operatorname{cl}_{X,p}:
 \operatorname{CH}^{p}(X)_{\mathbb Q}
 \twoheadrightarrow
 \operatorname{Hdg}^{p}(X)_{\mathbb Q}
\end{equation}
for every smooth complex projective variety \(X\) and every \(p\geq0\). This
is exactly the conventional rational formulation of the Hodge
statement. Here it appears as recognition of the chain
\(X_{\chi}\to\operatorname{Perf}\to\operatorname{CH}^{p}\to
\operatorname{Hdg}^{p}\), not as a condition used to construct it.

% Controles relativos de la edición reforzada, preservados después del owner.
\subsection{Relative auxiliary control: coend, kernels and based lifting}
\label{subsec:amp-hodge-presentacion-relativa}

\noindent\textbf{Status: NON\_GOVERNING\_CONTROL.}
The proprietary construction already proved is the APP--Mayer--Vietoris extension
with memory of \cref{prop:amp-hodge-extension-finita}, followed by the
Mittag--Leffler limit of
\cref{thm:amp-hodge-completitud-coinductiva}. The following kernels, cokernels and
coends are an additional relative presentation and local
falsifiers: they do not select \(\xi_\alpha\), are not premises of the proprietary
construction and cannot reopen
\eqref{eq:amp-hodge-evaluacion-alpha}.

To audit redundantly an extension already obtained by proposition
\ref{prop:amp-hodge-extension-finita}, this chart defines a relative operator
before fixing \(\alpha\). Set
\begin{equation}\label{eq:amp-hodge-nucleos-relativos}
 \mathscr K^{\mathrm S}_{N+1}
 :=\ker\rho_{N+1,N},
 \qquad
 \mathscr K^{\mathrm H}_{N+1}
 :=\ker q_{N+1,N},
\end{equation}
and let \(\Delta J_{N+1}\) be the ordered family of oriented cells incorporated
when passing from \(J_N\) to \(J_{N+1}\). For a cell
\(\nu=(Z,W,a,b)\), denote its perfect
realization by \(\mathcal C_\nu=\operatorname{Cone}(\kappa_{a,b})\).
The relative cellular module is
\begin{equation}\label{eq:amp-hodge-modulo-celular-relativo}
 \mathscr A_{N+1}^{\mathrm{rel}}
 :=H^0\!\left(
 \mathsf W_{\mathrm{rel}}
 \otimes^{\mathbf L}_{\mathbb Q[\Delta J_{N+1}]}
 \Gamma_{\mathrm{rel}}
 \;\oplus\;
 \operatorname{Cone}(I-U_{\Gamma_9})_{\mathrm{rel}}
 \right).
\end{equation}
In this expression, \(\Gamma_{\mathrm{rel}}(\nu)=\mathcal C_\nu\); the
coend imposes the Mayer--Vietoris and change-of-chart relations, and the
second summand preserves the terminal difference. The relation
\(\kappa_{b,a}\tau=-P\kappa_{a,b}\) fixes the signs of the opposite
generators, while \eqref{eq:hodge-suc-carry} fixes their composition.

There are two maps, both determined by the cellular generators,
\begin{equation}\label{eq:amp-hodge-mapas-relativos}
 \begin{aligned}
  a_{N+1}&:\mathscr A_{N+1}^{\mathrm{rel}}
    \longrightarrow\mathscr K^{\mathrm S}_{N+1},\\
  b_{N+1}&:\mathscr A_{N+1}^{\mathrm{rel}}
    \longrightarrow\mathscr K^{\mathrm H}_{N+1}.
 \end{aligned}
\end{equation}
The first inserts the cone as a state correction supported on the new
charts; the second publishes its incidence evaluation. By the local
compatibility of the localized character with Mayer--Vietoris, we have
\begin{equation}\label{eq:amp-hodge-factorizacion-relativa}
 e_{N+1}^{\mathrm{rel}}a_{N+1}=b_{N+1},
 \qquad
 e_{N+1}^{\mathrm{rel}}
 :=e_{N+1}^{\mathrm{Hdg}}\big|_{\mathscr K^{\mathrm S}_{N+1}}.
\end{equation}

The relative presentation makes it possible to isolate exactly the comparison whose
exactness is equivalent to the internal surjectivity of this redundant comparator.
It is neither a hypothesis nor a source of the proprietary surjectivity already
proved in \cref{prop:amp-hodge-extension-finita}. If
\(u:\nu\to\nu'\) runs through the arrows of
\(\Delta J_{N+1}\), set
\begin{equation}\label{eq:amp-hodge-relacion-coend}
 \mathscr T_{N+1}^{\mathrm{rel}}
 :=H^0\!\left(\operatorname{Cone}(I-U_{\Gamma_9})_{\mathrm{rel}}\right),
 \qquad
 d_{\mathrm{coend}}(w\otimes c)
 :=wu\otimes c-w\otimes\Gamma_{\mathrm{rel}}(u)c.
\end{equation}
The definition of the coend provides the modules
\begin{align}
 \mathscr R_{N+1}^{\mathrm{rel}}
 &:=
 \bigoplus_{u:\nu\to\nu'}
 \mathsf W_{\mathrm{rel}}(\nu')\otimes H^0(\mathcal C_\nu),
 \label{eq:amp-hodge-modulo-relaciones}\\
 \mathscr P_{N+1}^{\mathrm{rel}}
 &:=
 \left(
  \bigoplus_{\nu\in\Delta J_{N+1}}
  \mathsf W_{\mathrm{rel}}(\nu)\otimes H^0(\mathcal C_\nu)
 \right)\oplus\mathscr T_{N+1}^{\mathrm{rel}},
 \label{eq:amp-hodge-modulo-presentacion}
\end{align}
Before comparing it with the restriction kernel, the independent relative
codomain is formed
\begin{equation}\label{eq:amp-hodge-codominio-relativo-independiente}
 \mathscr O_{N+1}^{\mathrm{rel}}
 :=\operatorname{coker}
 \left(
  d_{\mathrm{coend}}:
  \mathscr R_{N+1}^{\mathrm{rel}}
  \longrightarrow\mathscr P_{N+1}^{\mathrm{rel}}
 \right),
 \qquad
 \pi_{N+1}^{\mathrm{rel}}:
 \mathscr P_{N+1}^{\mathrm{rel}}
 \twoheadrightarrow\mathscr O_{N+1}^{\mathrm{rel}}.
\end{equation}
This definition uses only the new cells, the Mayer--Vietoris and
change-of-chart relations and the terminal; it does not use
\(\mathscr K^{\mathrm H}_{N+1}\). Let
\(\widetilde b_{N+1}:\mathscr P_{N+1}^{\mathrm{rel}}\to
\mathscr K^{\mathrm H}_{N+1}\) be the map that publishes each cellular generator
in its new incidence and the terminal generator in the terminal difference.
Since \(\widetilde b_{N+1}d_{\mathrm{coend}}=0\), there is a unique map
\begin{equation}\label{eq:amp-hodge-comparacion-codominio-relativo}
 \chi_{N+1}^{\mathrm{rel}}:
 \mathscr O_{N+1}^{\mathrm{rel}}
 \longrightarrow\mathscr K^{\mathrm H}_{N+1},
 \qquad
 \chi_{N+1}^{\mathrm{rel}}\pi_{N+1}^{\mathrm{rel}}
 =\widetilde b_{N+1}.
\end{equation}

\begin{lemma}[Canonical comparison of the relative codomain]
\label{lem:amp-hodge-comparacion-codominio-relativo}
The independent cokernel has the exact presentation
\begin{equation}\label{eq:amp-hodge-presentacion-coend-relativa}
 \mathscr R_{N+1}^{\mathrm{rel}}
 \xrightarrow{\ d_{\mathrm{coend}}\ }
 \mathscr P_{N+1}^{\mathrm{rel}}
 \xrightarrow{\ \pi_{N+1}^{\mathrm{rel}}\ }
 \mathscr O_{N+1}^{\mathrm{rel}}\longrightarrow0 .
\end{equation}
If
\(\vartheta_{N+1}:\mathscr A_{N+1}^{\mathrm{rel}}
\xrightarrow{\sim}\mathscr O_{N+1}^{\mathrm{rel}}\)
is the canonical identification of the realization \(H^0\) of the coend with its
presentation, then
\begin{equation}\label{eq:amp-hodge-b-factor-cokernel}
 b_{N+1}
 =\chi_{N+1}^{\mathrm{rel}}\vartheta_{N+1}.
\end{equation}
The comparison with all new observations is measured by
\begin{equation}\label{eq:amp-hodge-obstrucciones-comparacion}
 \mathfrak o_{N+1}^{\ker}
 :=\ker\chi_{N+1}^{\mathrm{rel}},
 \qquad
 \mathfrak o_{N+1}^{\mathrm{coker}}
 :=\operatorname{coker}\chi_{N+1}^{\mathrm{rel}}.
\end{equation}
We have
\[
 \chi_{N+1}^{\mathrm{rel}}\text{ isomorphism}
 \quad\Longleftrightarrow\quad
 \mathfrak o_{N+1}^{\ker}
 =\mathfrak o_{N+1}^{\mathrm{coker}}=0.
\]
\end{lemma}

\begin{proof}
The relations
\(wu\otimes c-w\otimes\Gamma_{\mathrm{rel}}(u)c\) publish zero because
both terms are the same incidence written in two charts. Therefore
\eqref{eq:amp-hodge-comparacion-codominio-relativo} is well defined.
The exactness of \eqref{eq:amp-hodge-presentacion-coend-relativa} is the
universal property of the cokernel. The realization \(H^0\) of the coend
provides \(\vartheta_{N+1}\), and uniqueness of factorization through the
cokernel gives \eqref{eq:amp-hodge-b-factor-cokernel}. The last equivalence
is the elementary characterization of an isomorphism by the vanishing of its
kernel and cokernel. None of these steps proves that vanishing.
\end{proof}

Thus, the coend and the restriction kernel are constructed separately. In
this auxiliary chart, relative exactness is obtained if the specific
assertion
\(\mathfrak o_{N+1}^{\ker}
=\mathfrak o_{N+1}^{\mathrm{coker}}=0\) is proved. This obligation belongs only to
this comparator: it does not govern the proprietary
APP--Mayer--Vietoris extension or the coinductive completeness already proved.

\begin{lemma}[Cellular lifting under exact comparison]
\label{lem:amp-hodge-sobreyectividad-relativa}
Suppose
\(\mathfrak o_{N+1}^{\ker}
=\mathfrak o_{N+1}^{\mathrm{coker}}=0\).
The map \(b_{N+1}\) of
\eqref{eq:amp-hodge-mapas-relativos} is surjective. Moreover, the based
order of the charts determines a right inverse
\begin{equation}\label{eq:amp-hodge-sigma-relativa}
 \sigma_{N+1}^{\mathrm{rel}}:
 \mathscr K^{\mathrm H}_{N+1}\longrightarrow
 \mathscr A_{N+1}^{\mathrm{rel}},
 \qquad
 b_{N+1}\sigma_{N+1}^{\mathrm{rel}}=I.
\end{equation}
Consequently,
\begin{equation}\label{eq:amp-hodge-seccion-relativa}
 s_{N+1}^{\mathrm{rel}}
 :=a_{N+1}\sigma_{N+1}^{\mathrm{rel}}:
 \mathscr K^{\mathrm H}_{N+1}\longrightarrow
 \mathscr K^{\mathrm S}_{N+1}
 \quad\text{satisfies}\quad
 e_{N+1}^{\mathrm{rel}}s_{N+1}^{\mathrm{rel}}=I.
\end{equation}
\end{lemma}

\begin{proof}
An element of \(\mathscr K^{\mathrm H}_{N+1}\) vanishes when restricted to
\(J_N\); therefore, it is supported on the incidences of
\(\Delta J_{N+1}\). By the hypothesis on
\eqref{eq:amp-hodge-obstrucciones-comparacion},
\(\chi_{N+1}^{\mathrm{rel}}\) is an isomorphism. The exactness of
\eqref{eq:amp-hodge-presentacion-coend-relativa} then implies that its
restriction to each new
chart is a rational combination of the classes of the cones
\(\mathcal C_\nu\). On an intersection of charts, two such expressions
differ by the element
\(wu\otimes c-w\otimes\Gamma_{\mathrm{rel}}(u)c\) of
\eqref{eq:amp-hodge-relacion-coend}; passing to the coend makes that difference
zero. If the last chart completes a nonadic turn, the difference is not
eliminated: it remains in \(\mathscr T_{N+1}^{\mathrm{rel}}\). In this way every
relative class has a preimage under \(b_{N+1}\), proving
surjectivity.

To make the right inverse explicit, the generators are ordered by
depth, support, sheet and orientation, which are data of the atlas. Row
reduction of the matrix of \(b_{N+1}\) retains the first admissible generator
of each pivot column. If
\(\{\bar c_\mu\}_\mu\) is the resulting basis of
\(\mathscr K^{\mathrm H}_{N+1}\) and \(c_\mu\) the corresponding pivot
generator, one defines
\begin{equation}\label{eq:amp-hodge-sigma-coordenada}
 \sigma_{N+1}^{\mathrm{rel}}
 \left(\sum_\mu t_\mu\bar c_\mu\right)
 :=\sum_\mu t_\mu c_\mu.
\end{equation}
Then \(b_{N+1}(c_\mu)=\bar c_\mu\), from which one obtains
\eqref{eq:amp-hodge-sigma-relativa}.

The ordering used is the restriction of the global ordering of the atlas.
If \(r:J_{N+1}\to J'_{N+1}\) is a based refinement, let
\(r_{\mathscr A}\), \(r_{\mathrm H}\) and \(r_{\mathrm S}\) be the maps
induced on the cellular module, observations and relative states.
Refinement replaces a generator with its Mayer--Vietoris
expression and does not change the first global pivot of its class. Uniqueness
of the ordered normal form gives
\begin{equation}\label{eq:amp-hodge-naturalidad-seccion-relativa}
 r_{\mathscr A}\sigma_{N+1}^{\mathrm{rel}}
 =\sigma_{N+1}^{\prime\,\mathrm{rel}}r_{\mathrm H},
 \qquad
 r_{\mathrm S}s_{N+1}^{\mathrm{rel}}
 =s_{N+1}^{\prime\,\mathrm{rel}}r_{\mathrm H}.
\end{equation}
Thus, the right inverse is compatible with refinements, not merely with an
isolated presentation. The construction is carried out in oriented pairs;
therefore it respects
\(\kappa_{b,a}\tau=-P\kappa_{a,b}\). The coend relations impose the
carry cocycle, and the terminal generator is preserved instead of being reduced to
its visible phase. Finally, \eqref{eq:amp-hodge-factorizacion-relativa}
provides \eqref{eq:amp-hodge-seccion-relativa}. No step depends on
\(\alpha\) or on an algebraic class chosen beforehand.
\end{proof}

At depth zero, the based incidences simultaneously form the bases
of \(\mathscr S_0(X,p)\) and \(\mathscr H_0^p(X)\). If
\(\widehat c_\mu\) is the seed state publishing the observation
\(c_\mu\), the map
\begin{equation}\label{eq:amp-hodge-seccion-base}
 s_0^{\mathrm{base}}
 \left(\sum_\mu t_\mu c_\mu\right)
 :=\sum_\mu t_\mu\widehat c_\mu
 \quad\text{satisfies}\quad
 e_0^{\mathrm{Hdg}}s_0^{\mathrm{base}}=I.
\end{equation}
Likewise, the unitary degeneration of the atlas defines the prolongation with
zero relative coefficients
\begin{equation}\label{eq:amp-hodge-extension-cero}
 \iota_{N+1,N}:\mathscr S_N(X,p)\longrightarrow
 \mathscr S_{N+1}(X,p),
 \qquad
 \rho_{N+1,N}\iota_{N+1,N}=I.
\end{equation}
This prolongation inherits the terminal of \(s_N\); it does not reinitialize the history.

\subsection{Auxiliary commutativity check of the relative chart}
\label{subsec:amp-hodge-control-conmutatividad-relativa}

\noindent\textbf{Status: NON\_GOVERNING\_CONTROL.}
The computable defect of the relative chart is the transformation
\begin{equation}\label{eq:amp-hodge-defecto-puente-algebraico}
 \mathfrak d_N^{\mathrm{alg}}
 :=\operatorname{cl}_{X,p}\circ\operatorname{ch}^{\mathrm{loc}}_p
   \circ R_{\mathrm{Hdg},N}
   -\operatorname{pub}_N\circ e_N^{\mathrm{Hdg}}.
\end{equation}
The preceding proprietary construction already establishes that the diagram
\eqref{eq:amp-hodge-cuadrado-finito} commutes and, therefore,
\(\mathfrak d_N^{\mathrm{alg}}=0\). The preceding expression is a computable
falsifier and a redundant check of that identity, not a hypothesis
added to the theorem.

At the increment \(N\to N+1\), let
\[
 R_{\mathrm{Hdg},N+1}^{\mathrm{rel}}
 :=R_{\mathrm{Hdg},N+1}\big|_{\mathscr K^{\mathrm S}_{N+1}},
 \qquad
 \operatorname{pub}_{N+1}^{\mathrm{rel}}
 :=\operatorname{pub}_{N+1}\big|_{\mathscr K^{\mathrm H}_{N+1}}.
\]
The two maps with the codomains required by the localized character and
by the cycle class are
\begin{equation}\label{eq:amp-hodge-mapas-relativos-tipados}
 a_{N+1}^{\mathrm{mot}}
 :=R_{\mathrm{Hdg},N+1}^{\mathrm{rel}}a_{N+1},
 \qquad
 b_{N+1}^{\mathrm{coh}}
 :=\operatorname{pub}_{N+1}^{\mathrm{rel}}b_{N+1}.
\end{equation}
The well-typed relative defect is
\begin{equation}\label{eq:amp-hodge-cuadrado-relativo-tipado}
 \mathfrak d_{N+1}^{\mathrm{rel}}
 :=
 \operatorname{cl}_{X,p}\circ\operatorname{ch}^{\mathrm{loc}}_p
 \circ a_{N+1}^{\mathrm{mot}}
 -b_{N+1}^{\mathrm{coh}}.
\end{equation}
The relative publication identity checked by this auxiliary presentation is
\begin{equation}\label{eq:amp-hodge-identidad-algebraica-requerida}
 \boxed{\mathfrak d_{N+1}^{\mathrm{rel}}=0.}
\end{equation}
It is verified on the generators that
\(R_{\mathrm{Hdg},N+1}^{\mathrm{rel}}a_{N+1}\) realizes the corresponding perfect
cone, that the localized character publishes the same incidence
as \(b_{N+1}\), that both maps annihilate the Mayer--Vietoris
relations and preserve the same terminal difference. These
equalities certify the relative chart and do not replace, condition or
reopen the governing identity already established.

\subsection{Specialized validators of the Hodge chain}

\noindent\textbf{Status: NON\_GOVERNING\_CONTROL.}
The specialized checks of that chain include the models
\(C_3/K3\), the tritic curves and the Fermat sector through the \(405\)
planar incidences and their framework of \(486\) states; the Schur--Brauer
decomposition; semistable log-Perf descent; and, for every smooth
four-dimensional cubic \(X\), the universal line correspondence
\(P\subset F(X)\times X\). If
\(p:P\to X\), \(q:P\to F(X)\) and \(g\) is the polarization, the readers
\begin{equation}\label{eq:hodge-suc-fano}
 \Phi_X=p_*q^*,
 \qquad
 \Psi_X=q_*\bigl(p^*g^2\smile-\bigr),
 \qquad
 \Psi_X\Phi_X=-6I
\end{equation}
in the corresponding sector provide an explicit rational section;
the numerators are first realized in \(\operatorname{Perf}\).

In the Weil sector, for
\(A_m=E_i^m\times E_i^m\), take
\begin{equation}\label{eq:hodge-suc-weil-graficas}
 C^1=[\Gamma_1]-[X]-[Y],
 \qquad
 C^i=[\Gamma_i]-[X]-[Y],
 \qquad
 \operatorname{cl}(C^i+iC^1)=z\wedge\bar w,
\end{equation}
and
\begin{equation}\label{eq:hodge-suc-weil-determinante}
 P_m=\det_*(C^i+iC^1),
 \qquad
 \operatorname{cl}(P_m)
 =(-1)^{m(m-1)/2}m!\,w_+.
\end{equation}
The real and imaginary parts, normalized after constructing the
perfect numerators, realize the identity in the Weil sector. For
\(m=2\), the carry is two and does not require inverting three. The
Markman--Schoen result for complex abelian varieties of dimension four and
Weil type is recorded
as a recovered and typed external result; it is not attributed as an
original construction of this work.

\noindent\textbf{Hierarchy lock.}
The three control blocks of this section take as input identities
already proved by the owner. None of their comparators, vanishings or
specialized models is a premise of
\cref{thm:amp-hodge-completitud-coinductiva,thm:hodge-suc-generacion}; there is
no return causal edge from these controls to the owner.


\noindent\textbf{Status of the following control: NON\_GOVERNING\_CONTROL.}
The ablation receives the already proved Hodge conclusion and only falsifies the
suppression of the based antecedent; it contributes no premise to the owner and
cannot demote its theorems.

\begin{ablation}[Finite selector without history]
\label{abl:hodge-suc-eliptico}
Let \(E\) be a complex elliptic curve. There is no finite set of
rational representatives of the class of a point that is stable under all
translations of \(E\) and at the same time naturally selects a
representative without based data. This obstruction refutes the suppression of
rigidification; it does not affect the theorems
\ref{thm:hodge-suc-identidad} and \ref{thm:hodge-suc-generacion}.
\end{ablation}

\begin{proof}
If \(z\) has nonzero degree, the differences
\(t_x^{*}z-z\), as \(x\in E(\mathbb C)\) varies, range through a subgroup of \(\operatorname{Pic}^{0}(E)\)
of finite index up to isogeny via the Abel--Jacobi
map. Their rational span does not fit into the space
generated by a finite orbit. Therefore, a finite frame stable under all
translations cannot simultaneously preserve the representative and its
naturality. In \(X_{\chi}(E,1)\), by contrast, translation modifies the
history and its base even though the cohomological publication remains fixed. The
antecedent \(\xi_{\alpha}\) preserves precisely that datum, so the
hypothesis abolished by the control is not a premise of the theorem.
\end{proof}
\chapter{Coupled determinantal unit, rank--order of vanishing equality and Birch--Swinnerton--Dyer leading-term formula}
\chaptermark{Determinantal unit and Birch--Swinnerton--Dyer formula}

\noindent The specialization of the
\cref{thm:nucleo-continuidad-conexion-nonadica-conjunta} is
\[
 (\widehat X_\infty,\Gamma_9)
 \xrightarrow{\ \mathcal R_{\rm BSD}\ }
 P_E^{\rm gen}
 \longrightarrow(z_K,z_{\rm PT})
 \longrightarrow R_{\rm Boc}^{\rm der}
 \longrightarrow\text{rank and leading term of }L(E,s).
\]
The two publications start from the same based unit; the localization
sequence and Smith reduction precede the conclusion about
\(\Sha(E)\).

\subsection*{Genealogical object and Alexander--Whitney coalgebra}

Let \(E/\mathbb Q\) be an elliptic curve, and fix the coefficient system
\(A_n\), the depth \(m\) and the local and global orientations of the
Selmer complex. If
\(\operatorname{TPK}^{\mathrm{surv}}_{m}\) is the category of surviving
histories of the Topological Phase Kernel, define
\begin{equation}\label{eq:bsd-suc-gmn}
 G_{m,n}:=
 N_{*}\bigl(N\operatorname{TPK}^{\mathrm{surv}}_{m};A_n\bigr).
\end{equation}
For a nondegenerate simplex, the Alexander--Whitney diagonal is
\begin{equation}\label{eq:bsd-suc-aw}
 \Delta_{\mathrm{AW}}[x_0\to\cdots\to x_q]
 =\sum_{j=0}^{q}
 [x_0\to\cdots\to x_j]\otimes
 [x_j\to\cdots\to x_q].
\end{equation}
The counit equals one at vertices and zero in positive degrees. In
particular,
\begin{equation}\label{eq:bsd-suc-group-like}
 \Delta_{\mathrm{AW}}[x]=[x]\otimes[x],
 \qquad \varepsilon_{\mathrm{AW}}[x]=1.
\end{equation}
The group-like element property belongs to the complete-state vertex: it applies
before forgetting sheet, route, carry or memory.

Let \(P_E^{\mathrm{Man}}\to A[0]\) be the based Manin complex with its
augmentation, and let \(G_{m,n}\to A[0]\) be the augmentation induced by
\(\varepsilon_{\mathrm{AW}}\). The conservative target is the fiber
product
\begin{equation}\label{eq:bsd-suc-pullback}
 P_E^{\mathrm{gen}}
 :=P_E^{\mathrm{Man}}\times_{A[0]}G_{m,n}.
\end{equation}
If \(s:A[0]\to P_E^{\mathrm{Man}}\) is the based section, then
\begin{equation}\label{eq:bsd-suc-section}
 \widehat\Pi_0(c)
 :=\bigl(s\varepsilon_{\mathrm{AW}}(c),c\bigr)
\end{equation}
is a conservative section: its second projection returns \(c\), while
the first publishes its modular coordinate. The history and the modular shadow
are coupled in a single object.

\subsection*{The same unit in the Kato and Poitou--Tate channels}

The compatibilities of phase, sheet, corestriction and Hensel lifting
define on \(P_E^{\mathrm{gen}}\) a pair of publications
\begin{equation}\label{eq:bsd-suc-copub}
 P_E=(P_E^{K},P_E^{\mathrm{PT}}).
\end{equation}
Both components receive the same genealogical unit \(1_E\), so
that
\begin{equation}\label{eq:bsd-suc-zk-zpt}
 P_E^{K}(1_E)=z_K,
 \qquad P_E^{\mathrm{PT}}(1_E)=z_{\mathrm{PT}}.
\end{equation}
They are not two independent choices of generator.

Let \(L_{\mathrm{root}}=A_n z_{\mathrm{PT}}\) be the based root subcomplex,
normalized by
\(\lambda_{\mathrm{PT}}(z_{\mathrm{PT}})=1\), and let
\(p_{\mathrm{root}}\) be the chain-complex projector onto that subcomplex. Set
\(q_{\mathrm{root}}=I-p_{\mathrm{root}}\), which is also a chain-complex
projector; in particular, \(q_{\mathrm{root}}z_K\) is a cycle. Compatibility of the
counit in \eqref{eq:bsd-suc-group-like} with the pair
\eqref{eq:bsd-suc-copub} gives
\begin{equation}\label{eq:bsd-suc-normalizacion-raiz}
 \lambda_{\mathrm{PT}}(p_{\mathrm{root}}z_K)=1.
\end{equation}

Duality is formulated at the chain level. For a perfect complex
\(M\) over \(\Lambda\), let
\begin{equation}\label{eq:bsd-suc-d3}
 D_3(M):=
 \operatorname{RHom}_{\Lambda}(M^{\iota},\Lambda)[-3].
\end{equation}
If \(A\to C\to D_3(A)\) is the incidence block, define
\begin{equation}\label{eq:bsd-suc-orth-red}
 A^{\perp}:=\operatorname{Fib}\bigl(C\to D_3(A)\bigr),
 \qquad
 C^{\mathrm{red}}:=\operatorname{Cofib}\bigl(A\to A^{\perp}\bigr),
\end{equation}
and the reduced form induces
\begin{equation}\label{eq:bsd-suc-xorth}
 X^{\mathrm{orth}}:C^{\mathrm{red}}
 \xrightarrow{\ \sim\ }D_3(C^{\mathrm{red}}).
\end{equation}
On the reduced determinant line one obtains an involution
\(J_{\mathrm{red}}^2=I\) and compatible bases
\begin{equation}\label{eq:bsd-suc-jred}
 J_{\mathrm{red}}(\bar b_{\mathrm{PT}})=\bar b_K.
\end{equation}
The two sheets transport the same based scalar:
\begin{equation}\label{eq:bsd-suc-doble-hoja}
 u_{+}=u_{-}=\tau_{\mathrm{bas}}(D_K).
\end{equation}
The equality lives in the two lines identified by
\eqref{eq:bsd-suc-jred}; it is not replaced by an equation between a scalar and
its inverse. Likewise, \(D_3(D_K)=A^{\perp}\) belongs to the dual line and is not
confused with a second unoriented copy of the root line.

\begin{theorem}[Rigidity of the common root]
\label{thm:bsd-suc-raiz-filler}
With the preceding bases and orientations,
\begin{equation}\label{eq:bsd-suc-raiz}
 p_{\mathrm{root}}z_K=z_{\mathrm{PT}}.
\end{equation}
\end{theorem}

\begin{proof}
Write
\(p_{\mathrm{root}}z_K=u z_{\mathrm{PT}}\). Applying
\(\lambda_{\mathrm{PT}}\) and using the normalization of the base together with
\eqref{eq:bsd-suc-normalizacion-raiz} yields
\[
 1=\lambda_{\mathrm{PT}}(p_{\mathrm{root}}z_K)
  =u\lambda_{\mathrm{PT}}(z_{\mathrm{PT}})=u.
\]
This proves \eqref{eq:bsd-suc-raiz}; the identification
\eqref{eq:bsd-suc-jred} ensures that the conclusion is independent of the
sheet from which it is published.
\end{proof}

\begin{lemma}[Normal form of the complement]
\label{lem:bsd-suc-forma-normal}
In the finite--singular block with
free modules over \(A_n\), let \(g\) be primitive ---that is, part of a
basis of the target module--- and suppose that
\(\operatorname{im}\partial\) has no component in the direct summand
\(A_n r\). If the differential is given by
\begin{equation}\label{eq:bsd-suc-diferencial-normal}
 \partial f=3c e+b,
 \qquad \partial e=g,
 \qquad \partial b=-3c g,
\end{equation}
every cycle \(z_K=a r+u e+v b\) satisfies
\begin{equation}\label{eq:bsd-suc-ciclo-normal}
 u=3cv,
 \qquad
 z_{\mathrm{PT}}-z_K=(1-a)r-v\,\partial f,
 \quad z_{\mathrm{PT}}=r.
\end{equation}
The root equality \eqref{eq:bsd-suc-raiz} forces \(a=1\); therefore
\(h=-vf\) fills the complete difference.
\end{lemma}

\begin{proof}
The condition \(\partial z_K=0\) and
\eqref{eq:bsd-suc-diferencial-normal} give
\((u-3cv)g=0\). Since \(g\) is part of a basis of a free module,
comparison of its coefficient gives \(u=3cv\). Substituting this equality into
\(r-z_K\) produces \eqref{eq:bsd-suc-ciclo-normal}. The coefficient \(a\)
is precisely the coordinate of \(p_{\mathrm{root}}z_K\) in the basis
\(r=z_{\mathrm{PT}}\); theorem \ref{thm:bsd-suc-raiz-filler} gives
\(a=1\). Then
\(\partial(-vf)=-v\partial f=z_{\mathrm{PT}}-z_K\).
\end{proof}

The proprietary construction of the filling is therefore the explicit formula
\begin{equation}\label{eq:bsd-suc-hstar}
 h_*:=-vf,
 \qquad
 \partial h_*=z_{\mathrm{PT}}-z_K.
\end{equation}
No abstract homotopy is introduced as a premise: closedness fixes
\(u=3cv\), the root unit fixes \(a=1\) and the generator \(f\) directly
fills the complete difference.

\paragraph{Equivalent non-governing homotopic control.}
\noindent\textbf{Status: NON\_GOVERNING\_CONTROL.}
As a subsequent check, if the same complement is packaged by means of
a reader \(H_{\mathrm{FS}}\), the identity
\begin{equation}\label{eq:bsd-suc-hfs}
 \partial H_{\mathrm{FS}}+H_{\mathrm{FS}}\partial
 =q_{\mathrm{root}}
\end{equation}
reproduces \eqref{eq:bsd-suc-hstar} when applied to
\(q_{\mathrm{root}}z_K\). This reader is an auxiliary notation for the
contraction already performed by \(h_*=-vf\); it is not a prior obligation and
cannot reopen the explicit construction.

\subsection*{Derived Bockstein regulator of arbitrary order}

Let \(K\) be a field of characteristic zero,
\(\Lambda=K[[T]]\), and let
\begin{equation}\label{eq:bsd-suc-S}
 S(T):V_{\Lambda}\longrightarrow W_{\Lambda},
 \qquad
 V_{\Lambda}\simeq W_{\Lambda}\simeq\Lambda^{m},
\end{equation}
a based differential that becomes invertible over \(K((T))\). Fix
the line
\begin{equation}\label{eq:bsd-suc-lambda0}
 \lambda_0(S):=(\det_K V)^{-1}\otimes_K\det_K W,
 \quad
 V=V_{\Lambda}/TV_{\Lambda},\quad
 W=W_{\Lambda}/TW_{\Lambda}.
\end{equation}
A based Smith form has the expression
\begin{equation}\label{eq:bsd-suc-smith}
 U(T)S(T)V(T)
 =\operatorname{diag}
 \bigl(T^{a_1},\ldots,T^{a_s},1,\ldots,1\bigr),
 \qquad a_i\geq1,
\end{equation}
and
\begin{equation}\label{eq:bsd-suc-orden}
 d=\operatorname{ord}_{T}\det S(T)=\sum_{i=1}^{s}a_i.
\end{equation}

The \(T\)-adic filtration of the complex
\([V_{\Lambda}\xrightarrow{S}W_{\Lambda}]\) produces pages
\(V_q=E_q^0\), \(W_q=E_q^1\) and Bocksteins
\(\beta_q:V_q\to W_q\). Intrinsically,
\begin{equation}\label{eq:bsd-suc-pages}
 A_q:=V_q/V_{q+1},
 \qquad
 B_q:=\ker(W_q\twoheadrightarrow W_{q+1})
      =\operatorname{im}\beta_q,
\end{equation}
and \(\beta_q\) induces
\begin{equation}\label{eq:bsd-suc-beta-q}
 \bar\beta_q:A_q\xrightarrow{\ \sim\ }B_q,
 \qquad
 \tau_q:=\det(\bar\beta_q).
\end{equation}
The pages are finite and satisfy
\begin{equation}\label{eq:bsd-suc-carry-orden}
 d=\sum_{q\geq1}q\,\dim A_q
  =\sum_{q\geq1}\dim V_q.
\end{equation}
The first page sees
\(r_0=\dim V_1\); the higher-order carry is
\begin{equation}\label{eq:bsd-suc-carry-superior}
 c_I^{\deg}=d-r_0=\sum_{q\geq2}\dim V_q.
\end{equation}

The regular page is the induced isomorphism
\begin{equation}\label{eq:bsd-suc-pagina-regular}
 \bar S(0):A_0:=V/V_1
 \xrightarrow{\ \sim\ }
 B_0:=\operatorname{im}S(0),
 \qquad
 \tau_{\mathrm{reg}}:=\det\bar S(0).
\end{equation}
The exact sequences of \eqref{eq:bsd-suc-pages}, with their Koszul
signs, glue these trivializations into a single based element
\begin{equation}\label{eq:bsd-suc-rboc}
 R_{\mathrm{Boc}}^{\mathrm{der}}(S)
 :=\tau_{\mathrm{reg}}\otimes
   \bigotimes_{q\geq1}\tau_q
 \in\lambda_0(S),
\end{equation}
where the tensor symbol abbreviates the Knudsen--Mumford gluing
isomorphisms.

\begin{theorem}[Determinantal identity of the derived regulator]
\label{thm:bsd-suc-bockstein}
In the based line \(\lambda_0(S)\) one has
\begin{equation}\label{eq:bsd-suc-lt}
 R_{\mathrm{Boc}}^{\mathrm{der}}(S)
 =\operatorname{LT}_{T}(S)
 :=\left[T^{-d}\det S(T)\right]_{T=0}.
\end{equation}
The equality preserves the regular page, all positive pages and the
carry \eqref{eq:bsd-suc-carry-superior}.
\end{theorem}

\begin{proof}
In the Smith bases of \eqref{eq:bsd-suc-smith}, the regular page and each
positive page carry the canonical generator of their line. The Koszul
signs of gluing are exactly those needed to restore the order of
the total bases, so their product is the positive generator of
\(\lambda_0(S)\). On the other hand,
\[
 \det S(T)=\det U(T)^{-1}\det V(T)^{-1}T^d,
\]
and the initial term is
\(\det U(0)^{-1}\det V(0)^{-1}\). Upon undoing the change of bases, both
\eqref{eq:bsd-suc-rboc} and the initial term are multiplied by that same
factor. This proves \eqref{eq:bsd-suc-lt} without identifying scalars before
fixing the line and its basis.
\end{proof}

For each \(q\geq1\), the Poitou--Tate duality constructed on the same
page provides
\begin{equation}\label{eq:bsd-suc-jq}
 j_q:B_q\xrightarrow{\ \sim\ }
 A_q^{\vee}\otimes(I^q/I^{q+1}),
\end{equation}
and, after orienting \(I^q/I^{q+1}\), the pairing
\begin{equation}\label{eq:bsd-suc-altura-derivada}
 h_E^{(q)}(x,y)
 :=\bigl(j_q\bar\beta_q(x)\bigr)(y).
\end{equation}

% Ampliación sustantiva de los comparadores determinantes BSD.
% Módulo destinado al capítulo BSD de la Parte IV.

\section{Comparators of the common determinantal unit}
\label{sec:amp-bsd-comparadores}

The root equality \eqref{eq:bsd-suc-raiz} fixes a single unit in the Kato
and Poitou--Tate channels. From that unit, the comparators
transporting the derived regulator to its scalar publication
are now constructed. No comparator chooses a normalization after learning the
leading term: all arise from morphisms of complexes, exact
sequences and already oriented bases.

\subsection{The chain of determinant lines}
\label{subsec:amp-bsd-lineas}

For a based perfect complex \(C\), the convention used is
\begin{equation}\label{eq:amp-bsd-linea-determinante}
 \lambda(C):=\bigotimes_i
 \det H^i(C)^{(-1)^i}.
\end{equation}
Let \(\mathcal C_E^K\), \(\mathcal C_E^{\mathrm{Iw}}\),
\(\mathcal C_E^{\mathrm{PT}}\) and \(\mathcal C_E^{\mathrm{loc}}\) be,
respectively, the complex published by the Kato channel, its based
Iwasawa model, the self-dual reduced Poitou--Tate complex and the finite--singular
localization cone. Their lines are linked by
\begin{equation}\label{eq:amp-bsd-cadena-lineas}
 \mathcal L_E^K
 \xrightarrow{\ \mathfrak c_{K/\mathrm{FERS}}\ }
 \mathcal L_E^{\mathrm{Iw}}
 \xrightarrow{\ \mathfrak c_{\mathrm{PT}}\ }
 \mathcal L_E^{\mathrm{ht}}
 \xrightarrow{\ \mathfrak c_{\mathrm{STS}}\ }
 \mathcal L_E^{\mathrm{fin}}
 \xrightarrow{\ \mathfrak c_{\mathrm{tors}}\ }
 \mathcal L_E^{\mathrm{ar}}
 \xrightarrow{\ \mathfrak c_{\infty}\ }
 \mathcal L_E.
\end{equation}
In this formula, the abbreviation \(\mathrm{STS}\) denotes the joint step of
Smith, Tamagawa and Tate--Shafarevich. All the arrows of
\eqref{eq:amp-bsd-cadena-lineas} are isomorphisms of oriented lines.

\subsection{Kato--FERS chain comparison from the common unit}
\label{subsec:amp-bsd-comparacion-cadenas-kato-fers}

The equivalence giving rise to the first comparator can be constructed before
taking determinants. Let \(F^q\mathcal C\) be the filtration by nonadic
depth, identified with the \(T\)-adic filtration through the carry, and
write
\begin{equation}\label{eq:amp-bsd-complejos-kato-iwasawa}
 \mathcal C_E^K
 :=\operatorname{Tot}\!\left(
 P_E^K\widehat\Pi_0G_{m,n}\right),
 \qquad
 \mathcal C_E^{\mathrm{Iw}}
 :=[V_\Lambda\xrightarrow{S_E(T)}W_\Lambda].
\end{equation}
Let
\(\operatorname{car}_T:G_{m,n}\to\mathcal C_E^{\mathrm{Iw}}\) be the carry
reader assigning to a route \(\gamma\) the generator determined by its
initial and final faces, multiplied by
\(T^{\operatorname{Carr}(\gamma)}\). Multiplication of concatenations in
the Iwasawa model is denoted by \(\mu_{\mathrm{Iw}}\). The Alexander--Whitney
diagonal and the publication \(P_E^K\) then define, before
taking cohomology or determinants, the chain map
\begin{equation}\label{eq:amp-bsd-phi-kato-fers}
\begin{aligned}
 \Phi_{K/\mathrm{FERS}}&:
 \mathcal C_E^K\longrightarrow\mathcal C_E^{\mathrm{Iw}},\\
 \Phi_{K/\mathrm{FERS}}
 &:=
 \mu_{\mathrm{Iw}}\circ
 \bigl(P_E^K\widehat\otimes\operatorname{car}_T\bigr)
 \circ\Delta_{\mathrm{AW}},\\
 \Phi_{K/\mathrm{FERS}}(z_K)&=z_{\mathrm{Iw}}.
\end{aligned}
\end{equation}
where \(z_{\mathrm{Iw}}\) is the root generator induced by \(1_E\). Thus,
each normalized simplex is sent to the sum of its Alexander--Whitney
cuts, preserving in each summand the two faces and the carry
power. The formula does not consult the analytic leading term.

\begin{lemma}[Generator-by-generator criterion for equivalence]
\label{lem:amp-bsd-bases-emparejadas-kato-fers}
In each degree \(i\), let \(\mathcal B_i^K\) be the finite basis of normalized
simplices, ordered by depth, faces and carry, and let
\(\mathcal B_i^{\mathrm{Iw}}\) be the basis formed by the generators with the
same initial and final faces in the Iwasawa model. For a candidate
map
\[
 \upsilon_i:\mathcal B_i^K\longrightarrow
 \mathcal B_i^{\mathrm{Iw}},
\]
define the degree-zero residual
\begin{equation}\label{eq:amp-bsd-residual-generador-graduado}
 \varepsilon_i(b)
 :=\Phi_{K/\mathrm{FERS}}^i(b)\bmod T-\upsilon_i(b).
\end{equation}
If \(\upsilon_i\) is a bijection and
\(\varepsilon_i(b)=0\) for every \(b\in\mathcal B_i^K\), then the two
completed modules of degree \(i\) are free of the same rank
\(d_i=|\mathcal B_i^K|\) and, for each generator,
\begin{equation}\label{eq:amp-bsd-phi-en-cada-generador}
 \Phi_{K/\mathrm{FERS}}^i(b)
 =\upsilon_i(b)
  +T\sum_{c\in\mathcal B_i^{\mathrm{Iw}}}
    \lambda_{c,b}(T)c,
 \qquad \lambda_{c,b}(T)\in\Lambda.
\end{equation}
In particular,
\(\operatorname{gr}^{F}\Phi_{K/\mathrm{FERS}}^i=I\) on all
generators, not only on the root unit.
\end{lemma}

\begin{proof}
The vanishing of \eqref{eq:amp-bsd-residual-generador-graduado} says that the
column of \(b\) has constant term \(\upsilon_i(b)\). The remaining
coefficients belong to \(T\Lambda\), giving
\eqref{eq:amp-bsd-phi-en-cada-generador}. Bijectivity identifies the
bases and fixes the same rank. In those bases, the reduction modulo \(T\) of the
matrix of \(\Phi^i\) is the identity; hence
\(\operatorname{gr}^{F}\Phi^i=I\).
\end{proof}

The root equality
\(\Phi_{K/\mathrm{FERS}}(z_K)=z_{\mathrm{Iw}}\) fixes the common normalization.
The expression of \(P_E^K\) on each normalized simplex, the bijection
\(\upsilon_i\) and the residuals \(\varepsilon_i(b)\) provide a
generator-by-generator check of the equivalence already constructed; they neither choose nor
govern the BSD conclusion.

\begin{lemma}[Filtered equivalence and normalization of the root]
\label{lem:amp-bsd-equivalencia-kato-fers}
For the chain map constructed in
\eqref{eq:amp-bsd-phi-kato-fers}, the map \(\upsilon_i\) is the
basis correspondence induced by the faces and carry of each
simplex. The Alexander--Whitney formula verifies in each degree the
criterion of lemma \ref{lem:amp-bsd-bases-emparejadas-kato-fers}; this is a
check of the equivalence already constructed, not an additional premise.
The map \(\Phi_{K/\mathrm{FERS}}\) is a based equivalence of
perfect complexes. In the bases ordered by depth, its matrices in
each degree have the form
\begin{equation}\label{eq:amp-bsd-matriz-kato-fers}
 [\Phi_{K/\mathrm{FERS}}^i](T)=I+TB_i(T),
 \qquad B_i(T)\in M_{d_i}(\Lambda).
\end{equation}
Therefore, the unit
\begin{equation}\label{eq:amp-bsd-rho-determinantal}
 \rho_{E,\ell}(T)
 :=\prod_i\det\!\left(I+TB_i(T)\right)^{(-1)^i}
 \in\Lambda^\times
\end{equation}
satisface \(\rho_{E,\ell}(0)=1\).
\end{lemma}

\begin{proof}
The simplicial identity
\(\partial\Delta_{\mathrm{AW}}
=(\partial\otimes I+(-1)^{\deg}I\otimes\partial)
\Delta_{\mathrm{AW}}\), together with the compatibility of
\(P_E^K\) with faces and degeneracies, proves that
\eqref{eq:amp-bsd-phi-kato-fers} commutes with the differentials. On the associated
graded object, lemma
\ref{lem:amp-bsd-bases-emparejadas-kato-fers} identifies all the
generators and gives the identity. In those bases,
\eqref{eq:amp-bsd-phi-en-cada-generador} is precisely
\eqref{eq:amp-bsd-matriz-kato-fers}.

The completeness of \(\Lambda=K[[T]]\) makes the series converge
\begin{equation}\label{eq:amp-bsd-inversa-kato-fers}
 (I+TB_i(T))^{-1}
 =\sum_{n\geq0}(-TB_i(T))^n;
\end{equation}
these inverses also commute with the differentials, so they form the
based inverse of \(\Phi_{K/\mathrm{FERS}}\). Taking the alternating
determinant yields \eqref{eq:amp-bsd-rho-determinantal}; evaluating at
\(T=0\) gives one. Thus, the normalization is not postulated from the leading
term: it is the constant coordinate of an equivalence that is the identity
on the root unit.
\end{proof}

The first comparator is the determinant of
\(\Phi_{K/\mathrm{FERS}}\). If \(\mathfrak z_K\) denotes the element
coming from the unit \(1_E\), its image is characterized by
\begin{equation}\label{eq:amp-bsd-comparador-kato}
 \mathfrak c_{K/\mathrm{FERS}}(\mathfrak z_K)
 =\left[T^{-d}\rho_{E,\ell}(T)\det S_E(T)\right]_{T=0},
 \qquad
 \rho_{E,\ell}(0)=1.
\end{equation}
The unit of \(\rho_{E,\ell}\) is oriented and equals one at the root; it therefore does not
alter the order or the basis. Theorem
\ref{thm:bsd-suc-bockstein} transforms \eqref{eq:amp-bsd-comparador-kato} into
\begin{equation}\label{eq:amp-bsd-kato-bockstein}
 \mathfrak c_{K/\mathrm{FERS}}(\mathfrak z_K)
 =R_{\mathrm{Boc}}^{\mathrm{der}}(S_E)
 =\tau_{\mathrm{reg}}\otimes\bigotimes_{q\geq1}\tau_q.
\end{equation}

The Poitou--Tate comparator is constructed page by page. For
\(q\geq1\), the isomorphism \(j_q\) of
\eqref{eq:bsd-suc-jq} and the reduced Bockstein
\(\bar\beta_q\) induce
\begin{equation}\label{eq:amp-bsd-comparador-pt-q}
 \mathfrak c_{\mathrm{PT},q}(\tau_q)
 :=\det\!\left(j_q\circ\bar\beta_q\right)
 =\det h_E^{(q)}.
\end{equation}
The regular direction is transported through
\(\det\bar S_E(0)=\tau_{\mathrm{reg}}\). Knudsen--Mumford gluing and the
Koszul signs defined in \eqref{eq:bsd-suc-rboc} provide the unique
oriented isomorphism
\begin{equation}\label{eq:amp-bsd-comparador-pt}
 \mathfrak c_{\mathrm{PT}}
 \left(\tau_{\mathrm{reg}}\otimes\bigotimes_{q\geq1}\tau_q\right)
 =\tau_{\mathrm{reg}}^{\mathrm{ht}}
  \otimes\bigotimes_{q\geq1}\det h_E^{(q)}.
\end{equation}
On the first stable page, the form \(h_E^{(1)}\) is the
Poincaré--Néron--Tate form on the free Mordell--Weil lattice; its determinant is
\begin{equation}\label{eq:amp-bsd-regulador-nt}
 \det h_E^{(1)}=\operatorname{Reg}(E)
\end{equation}
in the basis induced by the root unit.

\begin{lemma}[Poitou--Tate unfolding of the arithmetic degree]
\label{lem:amp-bsd-grado-pt}
Let
\[
 \operatorname{Flag}_{\mathrm{Boc}}(E)
 :=\bigoplus_{q\geq1}V_q
\]
the finite flag of positive pages, each with the basis inherited from the
root unit. Orthogonal reduction and the isomorphisms \(j_q\) construct a
based isomorphism
\begin{equation}\label{eq:amp-bsd-psi-pt}
 \Psi_{\mathrm{PT}}:
 \bigl(E(\mathbb Q)/E(\mathbb Q)_{\mathrm{tors}}\bigr)\otimes K
 \xrightarrow{\ \sim\ }\operatorname{Flag}_{\mathrm{Boc}}(E).
\end{equation}
In particular,
\begin{equation}\label{eq:amp-bsd-rango-bandera}
 \operatorname{rank}E(\mathbb Q)
 =\sum_{q\geq1}\dim_KV_q
 =\operatorname{ord}_T\det S_E(T).
\end{equation}
\end{lemma}

\begin{proof}
The reduced complex is filtered by the powers of the augmentation ideal. In the
quotient at level \(q\), the regular direction has already been separated by
\(\bar S_E(0)\), and the only surviving block is \(V_q\). The equivalence
\(X^{\mathrm{orth}}\) identifies the opposite block with its dual, while
\(j_q\bar\beta_q\) provides the perfect pairing that lifts the
basis of the quotient to the preceding level. Starting at the last nonzero page and
repeating this lifting yields \(\Psi_{\mathrm{PT}}\).

In the bases ordered by depth, the matrix of the resulting map is
block triangular and all its diagonal blocks are identities: in the
first block because \(z_K\) and \(z_{\mathrm{PT}}\) have the same root, and
in the others because \(j_q\) and \(\bar\beta_q\) are based isomorphisms.
Therefore \(\Psi_{\mathrm{PT}}\) is invertible and loses no page.
Taking dimensions and using
\eqref{eq:bsd-suc-carry-orden} yields
\eqref{eq:amp-bsd-rango-bandera}. On taking determinants, the same
lifting is exactly the gluing of
\eqref{eq:amp-bsd-comparador-pt}; thus, equality of ranks and equality of
lines come from a single based map.
\end{proof}

\subsection{Localization triangle, Smith rank and finite factors}
\label{subsec:amp-bsd-smith-sha}

For each place \(v\), the finite local condition is a based subcomplex
\(\mathcal C_{E,v}^{\mathrm{fin}}\subset\mathcal C_{E,v}\), and one defines
\(\mathcal C_{E,v}^{\mathrm{sing}}\) by the triangle
\begin{equation}\label{eq:amp-bsd-triangulo-local}
 \mathcal C_{E,v}^{\mathrm{fin}}
 \xrightarrow{\ i_v\ }\mathcal C_{E,v}
 \longrightarrow\mathcal C_{E,v}^{\mathrm{sing}}
 \longrightarrow\mathcal C_{E,v}^{\mathrm{fin}}[1].
\end{equation}
If \(\operatorname{res}_v\) denotes the global restriction, the localization
map is, at chain level,
\begin{equation}\label{eq:amp-bsd-mapa-localizacion}
 \ell_E:
 \mathcal C_E^{\mathrm{glob,red}}\oplus
 \bigoplus_v\mathcal C_{E,v}^{\mathrm{fin}}
 \longrightarrow\bigoplus_v\mathcal C_{E,v},
 \qquad
 \ell_E\bigl(x,(y_v)_v\bigr)
 :=\bigl(\operatorname{res}_v x-i_vy_v\bigr)_v .
\end{equation}
The definition by cone of the reduced Selmer complex gives the triangle
\begin{equation}\label{eq:amp-bsd-triangulo-localizacion}
 \mathcal C_E^{\mathrm{red}}
 \longrightarrow
 \mathcal C_E^{\mathrm{glob,red}}\oplus
 \bigoplus_v\mathcal C_{E,v}^{\mathrm{fin}}
 \xrightarrow{\ \ell_E\ }
 \bigoplus_v\mathcal C_{E,v}
 \longrightarrow\mathcal C_E^{\mathrm{red}}[1].
\end{equation}
Therefore, local--global exactness comes from the cone of a map written in
\eqref{eq:amp-bsd-mapa-localizacion}; it is not subsequently added as a relation
between cardinalities.

In \eqref{eq:amp-bsd-triangulo-localizacion} the common root summand
is canceled, which is the same in both channels by
\eqref{eq:bsd-suc-raiz}. The resulting complement is denoted by
\begin{equation}\label{eq:amp-bsd-cono-local-reducido}
 \mathcal C_E^{\mathrm{loc,red}}
 :=q_{\mathrm{root}}\operatorname{Cone}(\ell_E)[-1].
\end{equation}
After separating the free Mordell--Weil lattice and its dual through
\(X^{\mathrm{orth}}\), a choice of the already oriented integral bases
represents this complex by
\begin{equation}\label{eq:amp-bsd-complejo-smith}
 \mathcal C_E^{\mathrm{loc,red}}
 \simeq[M_E\xrightarrow{\ D_E\ }N_E].
\end{equation}

\begin{lemma}[Rational acyclicity and full rank]
\label{lem:amp-bsd-aciclicidad-rango}
The complex of \eqref{eq:amp-bsd-complejo-smith} is acyclic after
tensoring with \(\mathbb Q\). In particular,
\[
 D_E\otimes\mathbb Q:M_E\otimes\mathbb Q
 \xrightarrow{\ \sim\ }N_E\otimes\mathbb Q
\]
is an isomorphism; \(M_E\) and \(N_E\) have the same rank and \(D_E\) has
full rank.
\end{lemma}

\begin{proof}
Let \(x\) be a cycle of the complex
\(\mathcal C_E^{\mathrm{loc,red}}\otimes\mathbb Q\). The reduction has
eliminated the summand \(p_{\mathrm{root}}\), hence
\(q_{\mathrm{root}}x=x\). In the oriented integral bases used in
\eqref{eq:amp-bsd-complejo-smith}, the finite--singular complement
decomposes into the normal blocks of
\cref{lem:bsd-suc-forma-normal}. In each block \(\lambda\), the root
coefficient is already zero and the cycle is written
\[
 x_{\lambda}=u_{\lambda}e_{\lambda}
                  +v_{\lambda}b_{\lambda},
 \qquad
 \partial x_{\lambda}
   =(u_{\lambda}-3c_{\lambda}v_{\lambda})g_{\lambda}=0.
\]
Since \(g_{\lambda}\) belongs to the free basis, we have
\(u_{\lambda}=3c_{\lambda}v_{\lambda}\), and the explicit formula for the
differential gives
\[
 x_{\lambda}
   =v_{\lambda}(3c_{\lambda}e_{\lambda}+b_{\lambda})
   =\partial(v_{\lambda}f_{\lambda}).
\]
The sum is finite in each degree, so
\(x=\partial\bigl(\sum_{\lambda}v_{\lambda}f_{\lambda}\bigr)\).
Every cycle is therefore a boundary by the same normal form that produced
the filler \eqref{eq:bsd-suc-hstar}, without introducing
\(H_{\mathrm{FS}}\) as a premise. The equality
\(p_{\mathrm{root}}z_K=z_{\mathrm{PT}}\) is what allows the root
direction to be canceled before this calculation; without it a degree-zero
class would remain. The self-dual equivalence
\(X^{\mathrm{orth}}:\mathcal C_E^{\mathrm{red}}\simeq
D_3(\mathcal C_E^{\mathrm{red}})\) pairs the two remaining degrees and
transports their bases with complementary opposite orientation. From this one
obtains \(\operatorname{rank}M_E=\operatorname{rank}N_E\).
Finally, acyclicity of a two-term complex of equal rank
is equivalent to \(D_E\otimes\mathbb Q\) being invertible.
\end{proof}

The lemma makes it possible to apply the Smith normal form without presupposing any of
its divisors. There exist oriented changes of basis \(U,V\) such that
\begin{equation}\label{eq:amp-bsd-morfismo-smith}
 U D_E V=\operatorname{diag}(s_1,\ldots,s_t),
 \qquad s_i\in\mathbb Z\setminus\{0\}.
\end{equation}
Therefore,
\begin{equation}\label{eq:amp-bsd-grupo-smith}
 F_E:=\operatorname{coker}D_E
 \simeq\bigoplus_{i=1}^{t}\mathbb Z/|s_i|\mathbb Z,
 \qquad
 |F_E|=\prod_{i=1}^{t}|s_i|<\infty.
\end{equation}

\begin{lemma}[Finite localization sequence]
\label{lem:amp-bsd-exactitud-sha-tamagawa}
The torsion fragment of the long cohomology sequence of
\eqref{eq:amp-bsd-triangulo-localizacion} is
\begin{equation}\label{eq:amp-bsd-exacta-sha-tamagawa}
 0\longrightarrow\operatorname{Sha}(E)
 \xrightarrow{\ \iota_{\mathrm{Sha}}\ }F_E
 \xrightarrow{\ \partial_{\mathrm{loc}}\ }
 \bigoplus_v\Phi_v(\mathbb F_v)
 \longrightarrow0,
 \qquad c_v=|\Phi_v(\mathbb F_v)|.
\end{equation}
\end{lemma}

\begin{proof}
A global cocycle represents a class of
\(\operatorname{Sha}(E)\) precisely when its restriction at each place
admits a primitive in
\(\mathcal C_{E,v}^{\mathrm{fin}}\). The class of the pair formed by that
cocycle and its local primitives in the cone of
\eqref{eq:amp-bsd-mapa-localizacion} defines
\(\iota_{\mathrm{Sha}}\). If this class is a boundary, the global component is a
boundary and the Tate--Shafarevich class is zero; thus,
\(\iota_{\mathrm{Sha}}\) is injective.

The map \(\partial_{\mathrm{loc}}\) takes a class of the cone, restricts it to
each quotient
\(\mathcal C_{E,v}^{\mathrm{sing}}\) of
\eqref{eq:amp-bsd-triangulo-local} and preserves its component class. The
identification
\[
 H^0(\mathcal C_{E,v}^{\mathrm{sing}})_{\mathrm{tors}}
 \simeq\Phi_v(\mathbb F_v)
\]
is induced by the finite--singular quotient, not by the order of the group.
A class of the cone has zero residue exactly when its local
representatives can be chosen in the finite subcomplexes; by the preceding
definition, that kernel is the image of \(\iota_{\mathrm{Sha}}\).
Finally, the last arrow of
\eqref{eq:amp-bsd-triangulo-local} lifts each component class to the
local complex. The exactness of the cone
\eqref{eq:amp-bsd-triangulo-localizacion}, after removing the two
free summands paired by \(X^{\mathrm{orth}}\), turns that
lifting into a preimage under \(\partial_{\mathrm{loc}}\). This proves
surjectivity and, consequently, all of
\eqref{eq:amp-bsd-exacta-sha-tamagawa}.
\end{proof}

Now, and only after lemma
\ref{lem:amp-bsd-aciclicidad-rango}, the group \(F_E\) is finite. The
injection of \eqref{eq:amp-bsd-exacta-sha-tamagawa} then proves the
finiteness of \(\operatorname{Sha}(E)\). Taking orders in that exact
sequence yields
\begin{equation}\label{eq:amp-bsd-cardinal-smith}
 |F_E|=|\operatorname{Sha}(E)|\prod_v c_v.
\end{equation}
This is the comparator
\(\mathfrak c_{\mathrm{STS}}\): it transports the Smith elementary factors
to the Tate--Shafarevich and Tamagawa factors while preserving their order in the line,
instead of inserting their cardinalities once the leading term has been calculated.

The torsion part is constructed by applying the determinant to the Mordell--Weil
lattice and its Pontryagin dual. The two finite factors are dual and
contribute
\begin{equation}\label{eq:amp-bsd-comparador-torsion}
 \mathfrak c_{\mathrm{tors}}:quad
 u\longmapsto
 \frac{u}{|E(\mathbb Q)_{\mathrm{tors}}|^2}.
\end{equation}
The Archimedean comparator is obtained by integrating the oriented Néron differential
over the based real cycle:
\begin{equation}\label{eq:amp-bsd-comparador-periodo}
 \Omega_E:=\int_{E(\mathbb R)}|\omega_E|,
 \qquad
 \mathfrak c_\infty:quad u\longmapsto\frac{u}{\Omega_E}.
\end{equation}
The differential and the cycle are fixed before evaluating \(L(E,s)\); therefore
\eqref{eq:amp-bsd-comparador-periodo} does not retrospectively normalize the
analytic section.

\subsection{Causal separation of primitives, identities and consequences}
\label{subsec:amp-bsd-separacion-causal}

The generator-by-generator criterion of
lemma \ref{lem:amp-bsd-bases-emparejadas-kato-fers} audits in each degree the
equivalence already constructed. It does not condition its existence. The logical order
of the package is
\begin{equation}\label{eq:amp-bsd-diagrama-causal}
\begin{array}{c}
 \begin{gathered}
  1_E,\ \Delta_{\mathrm{AW}},\ P_E^K,\ P_E^{\mathrm{PT}},
  \ h_*=-vf,\ X^{\mathrm{orth}},\\
  \{\,\mathcal C_{E,v}^{\mathrm{fin}}\to\mathcal C_{E,v}\,\}_v,\\
  \text{Mordell--Weil and Pontryagin bases;}\\
  \text{Néron differential and cycle}
 \end{gathered}
 \\[2mm]\Downarrow\\[-1mm]
 \begin{gathered}
  \Phi_{K/\mathrm{FERS}}\text{ equivalence},\quad
  \rho_{E,\ell}(0)=1,\quad
  p_{\mathrm{root}}z_K=z_{\mathrm{PT}},\\
  R_{\mathrm{Boc}}^{\mathrm{der}}=\operatorname{LT}_T(S_E),\quad
  \det(j_q\bar\beta_q)=\det h_E^{(q)},\\
  H^*(\mathcal C_E^{\mathrm{loc,red}}\otimes\mathbb Q)=0,\quad
  D_E\otimes\mathbb Q\text{ isomorphism},\\
  0\to\operatorname{Sha}(E)\to F_E
  \to\bigoplus_v\Phi_v(\mathbb F_v)\to0
 \end{gathered}
 \\[2mm]\Downarrow\\[-1mm]
 \begin{gathered}
  d_{\mathrm{an}}=\operatorname{rank}E(\mathbb Q),\qquad
  \operatorname{Reg}(E),\qquad
  |\operatorname{Sha}(E)|\prod_vc_v,\\
  |E(\mathbb Q)_{\mathrm{tors}}|^{-2},\qquad
  \Omega_E^{-1},\qquad
  \text{equality of the leading term}.
 \end{gathered}
\end{array}
\end{equation}
The first row contains only already based objects and maps, together with
the indicated generator-by-generator verification; the second, identities
proved by those maps; the third, their consequences on the determinant
line. In particular, no factor in the last row is
used to retroactively normalize an arrow in the first.

\begin{proposition}[Exact package of comparators]
\label{prop:amp-bsd-paquete-comparadores-exacto}
The comparators of
\eqref{eq:amp-bsd-cadena-lineas} come from a single package of maps of
complexes and satisfy, in the order of construction,
\begin{equation}\label{eq:amp-bsd-paquete-identidades}
 \begin{gathered}
 \Phi_{K/\mathrm{FERS}}(z_K)=z_{\mathrm{Iw}},\qquad
 \operatorname{gr}^{F}\Phi_{K/\mathrm{FERS}}=I,\\
 \det\nolimits_{\mathrm{alt}}\Phi_{K/\mathrm{FERS}}
   =\rho_{E,\ell}(T),\qquad \rho_{E,\ell}(0)=1,\\
 \det(j_q\bar\beta_q)=\det h_E^{(q)}\quad(q\geq1),\\
 \Psi_{\mathrm{PT}}:
  \bigl(E(\mathbb Q)/E(\mathbb Q)_{\mathrm{tors}}\bigr)\otimes K
  \xrightarrow{\sim}\operatorname{Flag}_{\mathrm{Boc}}(E),\\
 p_{\mathrm{root}}z_K=z_{\mathrm{PT}},\qquad
 h_*=-vf,\quad \partial h_*=z_{\mathrm{PT}}-z_K,\\
 H^*(\mathcal C_E^{\mathrm{loc,red}}\otimes\mathbb Q)=0,\qquad
 U D_E V=\operatorname{diag}(s_1,\ldots,s_t),\\
 0\longrightarrow\operatorname{Sha}(E)
  \longrightarrow F_E
  \longrightarrow\displaystyle\bigoplus_v\Phi_v(\mathbb F_v)
  \longrightarrow0 .
 \end{gathered}
\end{equation}
The induced composition on determinant lines is exactly
\begin{equation}\label{eq:amp-bsd-paquete-composicion}
 \det(\Phi_{K/\mathrm{FERS}})
 \xrightarrow{\ \det(j_q\bar\beta_q)\ }
 \det(D_E)
 \xrightarrow{\ \mathrm{Smith}\ }
 \frac{|\operatorname{Sha}(E)|\prod_vc_v}
      {|E(\mathbb Q)_{\mathrm{tors}}|^2\,\Omega_E},
\end{equation}
with the Knudsen--Mumford signs and the orientations of
\eqref{eq:amp-bsd-cadena-lineas}. In particular, the right-hand side of
\eqref{eq:amp-bsd-paquete-composicion} is the final coordinate of the transported
element, not data used to construct any of the arrows.
\end{proposition}

\begin{proof}
The first row of \eqref{eq:amp-bsd-paquete-identidades} is
\eqref{eq:amp-bsd-phi-kato-fers} and
\eqref{eq:amp-bsd-matriz-kato-fers}--
\eqref{eq:amp-bsd-rho-determinantal}. The second is
\eqref{eq:amp-bsd-comparador-pt-q} and lemma
\ref{lem:amp-bsd-grado-pt}. The third combines equality of the root with
the explicit filler \eqref{eq:bsd-suc-hstar}; its linear extension in the
normal form induces the auxiliary reader used in lemma
\ref{lem:amp-bsd-aciclicidad-rango}. On restricting to the complement of the
root, that reduction contracts every rational cycle. The last row is
\eqref{eq:amp-bsd-morfismo-smith} and
\eqref{eq:amp-bsd-exacta-sha-tamagawa}. Applying the determinant functor to
those identities, in that same order, gives
\eqref{eq:amp-bsd-paquete-composicion}. No equality is obtained
by first comparing the two final scalars.
\end{proof}

\begin{lemma}[Causal independence of the leading term]
\label{lem:amp-bsd-independencia-termino-principal}
Let \(\mathscr P_E\) be the data set
\begin{equation}\label{eq:amp-bsd-primitivas-paquete}
 \mathscr P_E:=
 \left(
  1_E,\Delta_{\mathrm{AW}},P_E^K,P_E^{\mathrm{PT}},
  h_*=-vf,X^{\mathrm{orth}},
  \{\mathcal C_{E,v}^{\mathrm{fin}}\xrightarrow{i_v}
     \mathcal C_{E,v}\}_v,
  \text{bases and orientations},
  \omega_E,\gamma_E^{\mathbb R}
 \right),
\end{equation}
where \(\omega_E\) is the Néron differential and
\(\gamma_E^{\mathbb R}\) the oriented real cycle. All maps of the
package in proposition
\ref{prop:amp-bsd-paquete-comparadores-exacto} are determined by
\(\mathscr P_E\). The following numerals are not part of \(\mathscr P_E\)
\[
 L^{(r)}(E,1),\quad r,\quad\operatorname{Reg}(E),\quad
 |\operatorname{Sha}(E)|,\quad(c_v)_v,\quad
 |E(\mathbb Q)_{\mathrm{tors}}|,\quad\Omega_E.
\]
Those values are obtained, respectively, by evaluating the analytic section,
taking the degree of the flag, calculating determinants and Smith forms,
taking orders of finite groups and integrating
\(\omega_E\) over \(\gamma_E^{\mathbb R}\).
\end{lemma}

\begin{proof}
The formula \eqref{eq:amp-bsd-phi-kato-fers} uses only
\(1_E,\Delta_{\mathrm{AW}},P_E^K\) and the carry. The Bockstein and
Poitou--Tate maps use the filtration, \(P_E^{\mathrm{PT}}\), the bases and
self-duality. The cone of \eqref{eq:amp-bsd-mapa-localizacion} uses the
maps \(i_v\) and the local restrictions; its contraction and Smith
form use the explicit filler \eqref{eq:bsd-suc-hstar}, its auxiliary linear
extension, \(X^{\mathrm{orth}}\) and the integral bases. The last
two comparators use the torsion lattices and the
pair \((\omega_E,\gamma_E^{\mathbb R})\). This is an exhaustive enumeration
of the arguments appearing in the construction formulas. The
numerals listed in the statement appear only after applying
cohomology, determinant, cardinality or integration, so they cannot
retrospectively select the maps.
\end{proof}

\subsection{Composition, degrees and leading term}
\label{subsec:amp-bsd-composicion}

Define the total comparator by composition in the written order:
\begin{equation}\label{eq:amp-bsd-comparador-total}
 \mathfrak C_E:=
 \mathfrak c_\infty\circ
 \mathfrak c_{\mathrm{tors}}\circ
 \mathfrak c_{\mathrm{STS}}\circ
 \mathfrak c_{\mathrm{PT}}\circ
 \mathfrak c_{K/\mathrm{FERS}}.
\end{equation}
Since \(p_{\mathrm{root}}z_K=z_{\mathrm{PT}}\), all arrows receive the
same based unit. The filler \(h_*=-vf\) of
\eqref{eq:bsd-suc-hstar} explicitly fills the difference; its linear
extension in the normal basis produces the auxiliary homotopic reader and proves that
no additional class modifies the element of the line during
transport. The reader is not a separate premise.

\begin{theorem}[Construction of the based comparators]
\label{thm:amp-bsd-comparadores-construidos}
For every elliptic curve \(E/\mathbb Q\) with the complexes, bases,
orientations and local data defined above, the maps
\eqref{eq:amp-bsd-phi-kato-fers},
\eqref{eq:amp-bsd-comparador-pt},
\eqref{eq:amp-bsd-triangulo-localizacion},
\eqref{eq:amp-bsd-comparador-torsion} and
\eqref{eq:amp-bsd-comparador-periodo}, together with lemmas
\ref{lem:amp-bsd-aciclicidad-rango} and
\ref{lem:amp-bsd-exactitud-sha-tamagawa}, construct the five comparators of
\eqref{eq:amp-bsd-cadena-lineas}. Their composition preserves the unit, degree
and orientation, and satisfies
\begin{equation}\label{eq:amp-bsd-igualdad-grados}
 d_{\mathrm{an}}=\operatorname{ord}_T\det S_E(T)
 =d_{\mathrm{ar}}=\operatorname{rank}E(\mathbb Q).
\end{equation}
Moreover, \(\operatorname{Sha}(E)\) is finite and the equality of the distinguished
element in \(\mathcal L_E\) is published as
\begin{equation}\label{eq:amp-bsd-termino-principal}
 \frac{L^{(r)}(E,1)}{r!\,\Omega_E}
 =\frac{|\operatorname{Sha}(E)|\,\operatorname{Reg}(E)
        \prod_v c_v}
       {|E(\mathbb Q)_{\mathrm{tors}}|^2},
 \qquad r=\operatorname{rank}E(\mathbb Q).
\end{equation}
\end{theorem}

\begin{proof}
The Kato--FERS comparator and the condition
\(\rho_{E,\ell}(0)=1\) identify the analytic order with
\(d=\operatorname{ord}_T\det S_E(T)\) and the initial term with
\(R_{\mathrm{Boc}}^{\mathrm{der}}(S_E)\). The identity
\eqref{eq:bsd-suc-lt} preserves the regular page and every positive page. The
isomorphisms \(j_q\) transport those pages to the derived heights.
Lemma \ref{lem:amp-bsd-grado-pt} constructs the isomorphism of the flag with the
free Mordell--Weil lattice and, together with
\eqref{eq:bsd-suc-carry-orden}, proves
\eqref{eq:amp-bsd-igualdad-grados}.

Lemma \ref{lem:amp-bsd-aciclicidad-rango} deduces rational acyclicity
directly from the explicit filler \eqref{eq:bsd-suc-hstar}, its linear
reduction of the complement, cancellation of the common root and
\(X^{\mathrm{orth}}\); from this the full rank of \(D_E\) is obtained, not
as an additional hypothesis. The Smith form
\eqref{eq:amp-bsd-morfismo-smith} then produces the finite group \(F_E\).
Lemma \ref{lem:amp-bsd-exactitud-sha-tamagawa} extracts from the triangles
\eqref{eq:amp-bsd-triangulo-local} and
\eqref{eq:amp-bsd-triangulo-localizacion} the finite exact sequence; only
afterwards are the finiteness of \(\operatorname{Sha}(E)\) and
\eqref{eq:amp-bsd-cardinal-smith} deduced. Finally,
\eqref{eq:amp-bsd-regulador-nt},
\eqref{eq:amp-bsd-comparador-torsion} and
\eqref{eq:amp-bsd-comparador-periodo} publish, respectively, the regulator,
the torsion quotient and the period. The root equality eliminates any relative
unit between the two publications. Evaluating the same based element at
both ends of \eqref{eq:amp-bsd-comparador-total} yields
\eqref{eq:amp-bsd-termino-principal}.
\end{proof}

\subsection{Premises and refutation criteria}
\label{subsec:amp-bsd-falsadores}

The construction starts from the genealogical unit \(1_E\) and the Alexander--Whitney
diagonal. It also uses the
coupled publications, all the
Bocksteins and the regular page, the explicit filler \(h_*=-vf\), the
reduced self-duality \(X^{\mathrm{orth}}\), the chain maps of the
localization triangle and the oriented bases of torsion and period. The
Kato--FERS equivalence and its normalization to one, the full rank of \(D_E\), the
exactness of \eqref{eq:amp-bsd-exacta-sha-tamagawa} and the finiteness of
\(\operatorname{Sha}(E)\) are conclusions of the preceding lemmas, not
independent premises. The generator-by-generator criterion of lemma
\ref{lem:amp-bsd-bases-emparejadas-kato-fers} is a nongoverning check
that locally audits the first comparator. The criteria that would refute a specific arrow
are:
\begin{enumerate}
 \item a root coefficient different from one, which would introduce a relative
       unit between \(z_K\) and \(z_{\mathrm{PT}}\);
 \item a noninvertible \(j_q\) or the omission of a Bockstein page, which
       would alter the degree or the height determinant;
 \item a zero Smith divisor, which would prevent the finiteness of \(F_E\);
 \item failure of exactness of
       \eqref{eq:amp-bsd-exacta-sha-tamagawa}, which would break the identification
       of the finite factors;
 \item an orientation reversal in torsion or period, which would change the
       based element even if it preserved its unoriented line.
\end{enumerate}
The memoryless Manin projection and visible rescaling modulo three
violate the first criterion; they therefore constitute necessity proofs for
those reduced models, not premises of theorem
\ref{thm:amp-bsd-comparadores-construidos}.


\begin{theorem}[Local--global determinantal comparison]
\label{thm:bsd-suc-local-global}
Let the based comparators already constructed in the
\cref{thm:amp-bsd-comparadores-construidos} be:
Kato--FERS/Iwasawa,
Poitou--Tate/Poincaré--Néron--Tate,
Smith--Tamagawa--Shafarevich, torsion and period.
If \(d_{\mathrm{an}}\) is the degree of the analytic section at the
central specialization and \(d_{\mathrm{ar}}\) the degree of the based Selmer
line, then
\begin{equation}\label{eq:bsd-suc-grados}
 d_{\mathrm{an}}=d=d_{\mathrm{ar}},
\end{equation}
and one obtains the equality of distinguished elements
\begin{equation}\label{eq:bsd-suc-linea-final}
 \mathfrak z_{\mathrm{an}}(E)=\mathfrak z_{\mathrm{ar}}(E)
 \quad\text{in}\quad\mathcal L_E.
\end{equation}
The analytic side is the publication of the initial term; the arithmetic
side is the publication of the complete regulator, the Smith
elementary factors,
the local factors, the torsion and the period.
\end{theorem}

\begin{proof}
The based construction gives
\([\Phi_{K/\mathrm{FERS}}^i](T)=I+TB_i(T)\) in each degree. The based
Kato--FERS comparator then expresses the corrected analytic section
as
\begin{equation}\label{eq:bsd-suc-fers}
 L_{\ell}^{\mathrm{corr}}(E,T)
 =\rho_{E,\ell}(T)\det S_E^{\mathrm{corr}}(T),
 \qquad \rho_{E,\ell}(0)=1.
\end{equation}
By \eqref{eq:bsd-suc-orden}, its central degree is \(d\); by
\eqref{eq:bsd-suc-lt}, its initial term is exactly
\(R_{\mathrm{Boc}}^{\mathrm{der}}\), not that regulator multiplied by an
undetermined unit. The isomorphisms \(j_q\) of
\eqref{eq:bsd-suc-jq} transport each \(\tau_q\) to the determinant of the
derived height of its page. The regular page transports the acyclic
summand that the positive pages do not see.

The Poincaré--Néron--Tate comparison identifies the degree of that height
line with \(d_{\mathrm{ar}}\). The Smith elementary factors decompose its
finite part into the Shafarevich and Tamagawa factors; the quotient of the integral
lattices supplies the square of the torsion, and the Archimedean reader supplies
the period. All these maps act on the same root line. By
theorem \ref{thm:bsd-suc-raiz-filler}, no relative unit remains between
the Kato and Poitou--Tate publications. The composition preserves the degree and
the based element, proving simultaneously
\eqref{eq:bsd-suc-grados} and \eqref{eq:bsd-suc-linea-final}. The corresponding scalar
formula is reserved for final recognition.
\end{proof}

\subsection*{Hypotheses and scope}

The BSD chain uses: the complete enriched state and its Alexander--Whitney
coalgebra; the fiber product \eqref{eq:bsd-suc-pullback}; the two
publications constructed on the same unit; the root line, its basis and its
counit; the perfect and saturated reduced quotient, together with
\(X^{\mathrm{orth}}\); the complete finite--singular complex; the finite
Bockstein flags, their Koszul signs and the regular page; and the
height, Smith, Tamagawa, Shafarevich, torsion and period comparisons
with a common orientation. The degree-by-degree verification of the paired
bases and the residuals \(\varepsilon_i(b)\) is recorded as a
local control of the realization, not as a governing premise of the closure.
Under this explicit starting structure,
\eqref{eq:bsd-suc-raiz}, \eqref{eq:bsd-suc-hstar},
\eqref{eq:bsd-suc-lt} and \eqref{eq:bsd-suc-linea-final} are exact
equalities.

Unit preservation, the two-sheet traversal and local--global
publication are formalized through the chain coalgebra, the genealogical
fiber product, the coupled publication, orthogonal reduction and the
Bockstein tower. The auxiliary identities separately verify the group-like element property,
preservation of history under \(\widehat\Pi_0\), the identity of the
unit in both channels, the reduced involution, the finite--singular normal
form and the determinantal identity of arbitrary order.

\paragraph{BSD ablation controls.}
\noindent\textbf{Status: NON\_GOVERNING\_CONTROL.}
The following three ablations receive the already proved identities and
falsify only reduced models. They contribute no premises to the owner,
do not condition its closure and have no causal return edge toward it.

\begin{ablation}[Manin projection without history]
\label{abl:bsd-suc-manin}
Replacing \(P_E^{\mathrm{gen}}\) with \(P_E^{\mathrm{Man}}\) erases data that
the coupled publication needs. Indeed, the route
\(\gamma=(\mathsf N\mathsf E)^{54}\) may have the same cell and the same
visible phase as the trivial route, but preserves distinct winding and memory
in \(G_{m,n}\). The Manin augmentation identifies them; the second
projection of \eqref{eq:bsd-suc-pullback} does not.
\end{ablation}

\begin{proof}
Every based lift to the Manin complex factors, up to homotopy, through the
section of the augmentation \(s\varepsilon_{\mathrm{AW}}\). It therefore depends only
on the visible coordinate. By contrast, the second component of
\(\widehat\Pi_0(c)\) is \(c\) itself, so it preserves the history
simplex and distinguishes the winding of \(\gamma\). The ablation refutes the
bare projection; it does not refute the fiber product used in the theorem.
\end{proof}

\begin{ablation}[Rescaling of the root visible only modulo three]
\label{abl:bsd-suc-twist}
In the specialization \(c=1\) of
\eqref{eq:bsd-suc-diferencial-normal}, the altered cycle
\begin{equation}\label{eq:bsd-suc-twist4}
 z_K^{(4)}=4r+3e+b
\end{equation}
satisfies
\begin{equation}\label{eq:bsd-suc-twist-defecto}
 z_{\mathrm{PT}}-z_K^{(4)}
 =\partial(-f)-3r.
\end{equation}
The defect \(-3r\) disappears modulo three and reappears modulo nine.
Consequently, \eqref{eq:bsd-suc-twist4} preserves a partial shadow, but not
the based unit or \eqref{eq:bsd-suc-raiz}.
\end{ablation}

\begin{proof}
The equality \(\partial f=3e+b\) immediately gives
\eqref{eq:bsd-suc-twist-defecto}. Modulo three, the second summand vanishes;
modulo nine it is a root class not divisible by a boundary of the complement.
Moreover, its root coordinate is four, whereas
\eqref{eq:bsd-suc-normalizacion-raiz} requires one. The example verifies the
necessity of the coupled pair and orientation; it introduces no exception
to the theorem that preserves both data.
\end{proof}

\begin{ablation}[Loss of regulator pages]
\label{abl:bsd-suc-bockstein}
The matrix \(\operatorname{diag}(T^2,T)\) shows that the first Bockstein does not
recover the higher-order carry. The matrix
\(\operatorname{diag}(2,T^2)\) has positive trivialization
\(\tau_2=1\), but
\(\tau_{\mathrm{reg}}=2=\operatorname{LT}_{T}(S)\). Therefore, neither the
first page nor the product of positive pages replaces the complete
regulator \eqref{eq:bsd-suc-rboc}.
\end{ablation}

\begin{proof}
In the first example, \(d=3\), \(r_0=2\) and
\(c_I^{\deg}=1\); page one does not contain that last degree of memory. In
the second, the only singular divisor is \(T^2\), but the regular direction
contributes the factor two. Omitting it would produce one instead of two. Both calculations
are diagonal cases of \eqref{eq:bsd-suc-smith} and verify the necessity of
all the pieces of \eqref{eq:bsd-suc-rboc}.
\end{proof}
\section*{Final mathematical recognition of Birch--Swinnerton--Dyer}

The preceding development concludes before employing its historical formulation
as a criterion.

In the domain declared by theorem
\ref{thm:bsd-suc-local-global}, the Smith comparison identifies the
finite component of the arithmetic line with the
Shafarevich--Tamagawa decomposition and proves the finiteness of
\(\operatorname{Sha}(E)\). Consequently, its order in the following scalar
publication is the cardinality of the identified finite group, not a formal
symbol added after the fact.

The scalar translation of \eqref{eq:bsd-suc-linea-final} provides, with
\(r=\operatorname{rank}E(\mathbb Q)\),
\begin{equation}\label{eq:bsd-suc-rango-final}
 \operatorname{ord}_{s=1}L(E,s)=r
\end{equation}
and
\begin{equation}\label{eq:bsd-suc-formula-final}
 \frac{L^{(r)}(E,1)}{r!\,\Omega_E}
 =\frac{
   \lvert\operatorname{Sha}(E)\rvert\,\operatorname{Reg}(E)\prod_v c_v}
  {\lvert E(\mathbb Q)_{\mathrm{tors}}\rvert^2}.
\end{equation}
This is the conventional formulation of Birch--Swinnerton--Dyer. The order,
the leading term and the arithmetic factors are not defined through one
another: they are coordinated publications of the based equality
\eqref{eq:bsd-suc-linea-final}, whose common root was previously fixed by
\eqref{eq:bsd-suc-raiz}.
